\chapter{12.11-4 等边三角形探究}

\section*{p001}
\subsection*{四、未分类题目（共 1 小题）}

\solutionlist[18]
\item
  已知: 等边 $\triangle ABC$ 的边长为 $a$ .
探究（1）：如图\circled{1}，过等边 $\triangle ABC$ 的顶点 $A$、 $B$、 $C$ 依次作 $AB$、 $BC$、 $CA$ 的垂线围成 $\triangle MNG$，求证： $\triangle MNG$ 是等边三角形且 $MN=\sqrt{3} a$
探究（2）：在等边 $\triangle ABC$ 内取一点 $O$，过点 $O$ 分别作 $OD \perp AB$、 $OE \perp BC$、 $OF \perp CA$，垂足分别为点 $D$、 $E$、 $F$ .
\circled{1}如图\circled{2}，若点 $O$ 是 $\triangle ABC$ 的重心，我们可利用三角形面积公式及等边三角形性质得到两个正确结论（不必证明）：结论1. $OD+OE+OF=\frac{\sqrt{3}}{2} a$；结论2. $AD+BE+CF=\frac{3}{2} a$；
\circled{2}如图\circled{3}，若点 $O$ 是等边 $\triangle ABC$ 内任意一点，则上述结论1，2是否仍然成立？如果成立，请给予证明；如果不成立，请说明理由.

  \begin{questionimagegroup}
  \questionimageitem{0.28\textwidth}{chuyi_xia_2-p001-cleaned-1.jpg}{(图\circled{1})}
  \questionimagegap
  \questionimageitem{0.28\textwidth}{chuyi_xia_2-p001-cleaned-2.jpg}{(图\circled{2})}
  \questionimagegap
  \questionimageitem{0.28\textwidth}{chuyi_xia_2-p001-cleaned-3.jpg}{(图\circled{3})}
  \questionimagegap
  \questionimageitem{0.28\textwidth}{chuyi_xia_2-p001-cleaned-4.jpg}{(图\circled{4})}
  \end{questionimagegroup}

\end{enumerate}

\section*{p002}
\subsection*{一、选择题（共 7 小题）}

\questionlist
\item
  \begin{questionwithimage}{chuyi_xia_2-p002-cleaned-1.jpg}
  如图，在一条公路 $CD$ 的同一侧有 $A$、 $B$ 两个村庄， $A$、 $B$ 与公路的距离 $AC$、 $BD$ 分别为 $500\mathrm{m}$ 和 $700\mathrm{m}$，且 $C$、 $D$ 两地相距 $500\mathrm{m}$，若要公路旁（在 $CD$ 上）建一个车站，则 $A$、 $B$ 两村庄到车站的距离之和最短是（\quad）
  \optionlist
  \item $1000\mathrm{m}$
  \item $1200\mathrm{m}$
  \item $1300\mathrm{m}$
  \item $1700\mathrm{m}$
  \end{enumerate}
  \end{questionwithimage}

\item
  \begin{questionwithimage}{chuyi_xia_2-p002-cleaned-2.jpg}
  如图，在等腰三角形 $ABC$ 中， $\angle ABC=120^\circ$，点 $P$ 是底边 $AC$ 上的一个动点， $M$、 $N$ 分别是 $AB$、 $BC$ 的中点，若 $PM+PN$ 的最小值为 6，则 $\triangle ABC$ 的周长为（\quad）
  \optionlist
  \item 12
  \item $12+6\sqrt{3}$
  \item $12+7\sqrt{3}$
  \item $12+10\sqrt{3}$
  \end{enumerate}
  \end{questionwithimage}

\item
  \begin{questionwithimage}{chuyi_xia_2-p002-cleaned-3.jpg}
  如图，正方形 $ABCD$ 的面积为 12， $\triangle ABE$ 是等边三角形，点 $E$ 在正方形 $ABCD$ 内，在对角线 $AC$ 上有一点 $P$，使 $PD+PE$ 最小，则这个最小值为（\quad）
  \optionlist
  \item $\sqrt{3}$
  \item $2\sqrt{3}$
  \item $2\sqrt{6}$
  \item $\sqrt{6}$
  \end{enumerate}
  \end{questionwithimage}

\item
  \begin{questionwithimage}{chuyi_xia_2-p002-cleaned-4.jpg}
  如图，四边形 ABCD 中， $\angle C=50^{\circ}$， $\angle B=\angle D=90^{\circ}$，（\quad）
  \optionlist
  \item F 分别是 B
  \item DC 上的点，当 $\triangle AEF$ 的周长最小时， $\angle EAF$ 的度数为（）
  \item $50^{\circ}$
  \item $60^{\circ}$
  \item $70^{\circ}$
  \item $80^{\circ}$
  \end{enumerate}
  \end{questionwithimage}

\item
  如图，点 $P$ 是 $\angle AOB$ 内任意一点， $OP=5cm$，点 $M$ 和点 $N$ 分别是射线 $OA$ 和射线 $OB$ 上的动点， $\triangle PMN$ 周长的最小值是 $5cm$，则 $\angle AOB$ 的度数是（\quad）
  \optionlist
  \item $25^{\circ}$
  \item $30^{\circ}$
  \item $35^{\circ}$
  \item $40^{\circ}$
  \end{enumerate}

\item
  \begin{questionwithimage}{chuyi_xia_2-p002-cleaned-5.jpg}
  如图，在平面直角坐标系中，A 点为直线 y=x 上一点，过 A 点作 $AB \perp x$ 轴于 B 点，若 OB=4，E 是 $OB$边上的一点，且 OE=3，点$P$为线段$AO$上的动点，则 $\triangle BEP$ 周长的最小值为（\quad）.
  \optionlist
  \item $4+2\sqrt{2}$
  \item $4+\sqrt{10}$
  \item 6
  \item $4\sqrt{2}$
  \end{enumerate}
  \end{questionwithimage}

\item
  \begin{questionwithimage}{chuyi_xia_2-p002-cleaned-6.jpg}
  如图，已知等边 $\triangle ABC$ 的边长为 6，点 $D$ 为 $AC$ 的中点，点 $E$ 为 $BC$ 的中点，点 $P$ 为 $BD$ 上一点，则 $PE+PC$ 的最小值为（\quad）
  \optionlist
  \item 3
  \item $3\sqrt{2}$
  \item $2\sqrt{3}$
  \item $3\sqrt{3}$
  \end{enumerate}
  \end{questionwithimage}

\end{enumerate}

\section*{p003}
\subsection*{四、未分类题目（共 2 小题）}

\solutionlist[8]
\item
  如图，在边长为1正方形ABCD中，（\quad）
  \optionlist
  \item 
  \item 
  \item H分别是A
  \item B
  \item C
  \item DA上的点，3AE=EB，有一只蚂蚁从E点出发，经过
  \item 
  \item H，最后回点$E$点，则蚂蚁所走的最小路程是（）
  \item 2
  \item 4
  \item $2\sqrt{2}$
  \item $3\sqrt{2}$
  \end{enumerate}

\item
  如图，在平面直角坐标系中，有 $A$(1,2)，$B$(3,3) 两点，现另取一点 $C$(a,1)，当 a=（\quad）时，AC+BC 的值最小.
  \optionlist
  \item 2
  \item $\frac{5}{3}$
  \item $\frac{11}{4}$
  \item 3
  \end{enumerate}

\end{enumerate}

\subsection*{二、填空题（共 5 小题）}

\blanklist[10]
\item
  如图， $\angle AOB=30^{\circ}$，点 $M$、 $N$ 分别在边 $OA$、 $OB$ 上，且 $OM=1$， $ON=3$。点 $P$、 $Q$ 分别在边 $OB$、 $OA$ 上，则 $MP+PQ+ON$ 的最小值是 \_\_\_\_.\underline{\hspace{2cm}}

\item
  如图，在直角坐标系中有四个点 $A(-6, 3)$ , $B(-2, 5)$ , $C(0, m)$ , $D(n, 0)$ , 当四边形 $ABCD$ 周长最短时, 则 $m=\_$ , $n=\_$ .\underline{\hspace{2cm}}

\item
  如图：点 $P$ 是 $\angle AOB$ 内一定点，点 $M$、 $N$ 分别在边 $OA$、 $OB$ 上运动，若 $\angle AOB=45^{\circ}$， $OP=3\sqrt{2}$，则 $\triangle PMN$ 的周长的最小值为 \_\_\_\_.\underline{\hspace{2cm}}

\item
  点 $A$、 $B$ 均在由面积为 1 的相同小矩形组成的网格的格点上，建立平面直角坐标系如图所示。若 $P$ 是 $x$ 轴上使得 $|PA-PB|$ 的值最大的点， $Q$ 是 $y$ 轴上使得 $QA+QB$ 的值最小的点，则 $OP \cdot OQ=$ \_\_\_\_.\underline{\hspace{2cm}}

\item
  若 $x > 0, y > 0$，且 $x+y=12$。则 $\sqrt{x^2+4}+\sqrt{y^2+9}$ 的最小值是 \_\_\_\_.\underline{\hspace{2cm}}

\end{enumerate}

\section*{p004}
\subsection*{四、未分类题目（共 3 小题）}

\solutionlist[15]
\item
  如图，已知 $\angle AOB$ 和 $C$、 $D$ 两点，求作一点 $P$ 使 $P$ 到 $\angle AOB$ 两边的距离相等，且使 $P$ 到 $C$ 两点的距离和最小.

  \centerquestionimage{0.35\textwidth}{chuyi_xia_2-p004-cleaned-1.jpg}

\item
  如图 $\triangle ABC$ 是边长为2的等边三角形， $D$ 是 $AB$ 边的中点， $P$ 是 $BC$ 边上的动点， $Q$ 是 $AC$ 边上的动点，当 $P$、 $Q$ 的位置在何处时，才能使 $\triangle DPQ$ 的周长最小？并求出这个最值.

  \centerquestionimage{0.35\textwidth}{chuyi_xia_2-p004-cleaned-2.jpg}

\item
  如图，在长方形 $ABCD$ 中， $O$ 为对角线 $AC$ 的中点， $P$ 是 $AB$ 上任意一点， $Q$ 是 $OC$ 上任意一点，已知： $AC=2, BC=1$ .
（1）求折线 OPQB 的长的最小值；
(2) 当折线 OPQB 的长最小时，试确定 Q 的位置.

  \centerquestionimage{0.35\textwidth}{chuyi_xia_2-p004-cleaned-3.jpg}

\end{enumerate}

\section*{p005}
\subsection*{四、未分类题目（共 1 小题）}

\solutionlist[18]
\item
  在一平直河岸 $l$ 同侧有 $A, B$ 两个村庄， $A, B$ 到 $l$ 的距离分别是 $3km$ 和 $2km, AB=akm(a > 1)$ . 现计划在河岸 $l$ 上建一抽水站 $P$，用输水管向两个村庄供水.
某班数学兴趣小组设计了两种铺设管道方案：图\circled{1}是方案一的示意图，设该方案中管道长度为 $d_{1}$ 且 $d_{1}=PB+BA(km)$ （其中 $BP \perp l$ 于点 $P$）；图\circled{2}是方案二的示意图，设该方案中管道长度为 $d_{2}$，且 $d_{2}$ $=PA+PB(km)$ （其中点 $A^{\prime}$ 与点 $A$ 关于 $l$ 对称， $A^{\prime}B$ 与 $l$ 交于点 $P$）.
观察计算：（1）在方案一中， $d_{1}=$ \_\_\_\_km（用含 $\sigma$ 的式子表示）；
（2）在方案一中，组长小宇为了计算 $d_{2}$ 的长，作了如图\circled{3}所示的辅助线，请你按小宇同学的思路计算， $d_{2}=$ \_\_\_\_ km（用含 $a$ 的式子表示）.
探索归纳：（1）\circled{1}当 a=4 时，比较大小： $d_{1}$ \_\_\_\_ $d_{2}$ （填“＞”、“＝”或“＜”）；
\circled{2}当a=6时，比较大小： $d_{1}$ \_\_\_\_ $d_{2}$ （填“>”、“=”或“<”）；
（2）请你参考方法指导，就 $a$ （当 $a > 1$ 时）的所有取值情况进行分析，要使铺设的管道长度较短，应选择方案一还是方案二？

  \begin{questionimagegroup}
  \questionimageitem{0.28\textwidth}{chuyi_xia_2-p005-cleaned-1.jpg}{图\circled{1}}
  \questionimagegap
  \questionimageitem{0.28\textwidth}{chuyi_xia_2-p005-cleaned-2.jpg}{图\circled{2}}
  \questionimagegap
  \questionimageitem{0.28\textwidth}{chuyi_xia_2-p005-cleaned-3.jpg}{图\circled{3}}
  \end{questionimagegroup}

\end{enumerate}

\section*{p006}
\subsection*{一、选择题（共 6 小题）}

\questionlist
\item
  \begin{questionwithimage}{chuyi_xia_2-p006-cleaned-1.jpg}
  如图，将 $\triangle ABC$ 绕着点 $C$ 顺时针旋转 $50^{\circ}$ 后得到 $\triangle A'B'C'$。若 $\angle A=40^{\circ}$ $\angle B'=110^{\circ}$，则 $\angle BCA'$ 的度数是（\quad）
  \optionlist
  \item $110^{\circ}$
  \item $80^{\circ}$
  \item $40^{\circ}$
  \item $30^{\circ}$
  \end{enumerate}
  \end{questionwithimage}

\item
  \begin{questionwithimage}{chuyi_xia_2-p006-cleaned-2.jpg}
  如图，在 $\triangle ABC$ 中， $\angle CAB=65^{\circ}$，将 $\triangle ABC$ 在平面内绕点 $A$ 旋转到 $\triangle AB'C'$ 的位置，使 $CC' \parallel AB$，则旋转角的度数为（\quad）
  \optionlist
  \item $35^{\circ}$
  \item $40^{\circ}$
  \item $50^{\circ}$
  \item $65^{\circ}$
  \end{enumerate}
  \end{questionwithimage}

\item
  \begin{questionwithimage}{chuyi_xia_2-p006-cleaned-3.jpg}
  如图，将 Rt$\triangle ABC$（其中 $\angle B=30^\circ$，$\angle C=90^\circ$）绕点 $A$ 按顺时针方向旋转到 $\triangle AB_{1}C_{1}$ 的位置，使得点 $C$、（\quad）
  \optionlist
  \item $B_{1}$ 在同一条直线上，那么旋转角等于（）
  \item 115$^\circ$
  \item 120$^\circ$
  \item 125$^\circ$
  \item 145$^\circ$
  \end{enumerate}
  \end{questionwithimage}

\item
  如图， $\triangle ABC$ 为等腰直角三角形， $\angle ACB=90^\circ$，将 $\triangle ABC$ 绕点 $A$ 逆时针旋转 $75^\circ$，得到 $\triangle AB'C'$，过点 $B'$ 作 $B'D \perp CA$，交 $CA$ 的延长线于点 $D$，若 $AC=6$，则 $AD$ 的长为（\quad）
  \optionlist
  \item 2
  \item 3
  \item $2\sqrt{3}$
  \item $3\sqrt{2}$
  \end{enumerate}

\item
  如图，在平面直角坐标系 $xOy$ 中， $\triangle A'B'C'$ 由 $\triangle ABC$ 绕点 $P$ 旋转得到，则点 $P$ 的坐标为（\quad）
  \optionlist
  \item (0, 1)
  \item (1, -1)
  \item (0, -1)
  \item (1, 0)
  \end{enumerate}

\item
  \begin{questionwithimage}{chuyi_xia_2-p006-cleaned-4.jpg}
  一副三角板叠在一起如图放置，最小锐角的顶点 $D$ 恰好放在等腰直角三角形的斜边上， $AC$ 与 $DM$， $DN$ 分别交于点 $E$， $F$，把 $\triangle DEF$ 绕点 $D$ 旋转到一定位置，使得 $DE=DF$，则 $\angle BDN$ 的度数是（\quad）
  \optionlist
  \item $105^{\circ}$
  \item $115^{\circ}$
  \item $120^{\circ}$
  \item $135^{\circ}$
  \end{enumerate}
  \end{questionwithimage}

\end{enumerate}

\section*{p007}
\subsection*{四、未分类题目（共 2 小题）}

\solutionlist[7]
\item
  \begin{questionwithimage}{chuyi_xia_2-p007-cleaned-1.jpg}
  如图，在平面直角坐标系 $xOy$ 中， $OA$ 与 $x$ 轴正方向的夹角是 $60^\circ$，作 $AB \perp x$ 轴于点 $B$，将 $\triangle ABO$ 绕点 $B$ 逆时针旋转 $60^\circ$ 得到 $\triangle CBD$。若点 $B$ 的坐标为 $(2, 0)$，则点 $C$ 的坐标为（\quad）
  \optionlist
  \item $(-1, \sqrt{3})$
  \item $(-2, \sqrt{3})$
  \item $(-\sqrt{3}, 1)$
  \item $(-\sqrt{3}, 2)$
  \end{enumerate}
  \end{questionwithimage}

\item
  \begin{questionwithimage}{chuyi_xia_2-p007-cleaned-2.jpg}
  如图，在 $\triangle ABC$ 中， $AB=AC$， $\angle BAC=90^\circ$，直角 $\angle EPF$ 的顶点 $P$ 是 $BC$ 的中点，两边 $PE$， $PF$ 分别交 $AB$， $AC$ 于点 $E$， $F$，现给出以下四个结论：（1） $AE=CF$；（2） $\triangle EPF$ 是等腰直角三角形；（3） $S_{\text{四边形} AEPF}=\frac{1}{2} S_{\triangle ABC}$；
（4）当$\angle$EPF在$\triangle ABC$内绕顶点$P$旋转时始终有EF=AP.（点$E$不与（\quad）
  \optionlist
  \item B重合），上述结论中是正确的结论的有（）
  \item 1个
  \item 2个
  \item 3个
  \end{enumerate}
  \end{questionwithimage}

\end{enumerate}

\subsection*{二、填空题（共 5 小题）}

\blanklist[9]
\item
  在平面直角坐标系中，Rt$\triangle OAB$的顶点$A$ 的坐标为 $\left(\sqrt{3.1}\right)$，若将$\triangle OAB$绕 O 点，逆时针旋转 $60^{\circ}$ 后，B 点到达 B' 点，则点$B$' 的坐标是 \_\_\_\_。\underline{\hspace{2cm}}

\item
  如图，在 $\triangle ABC$ 中， $\angle A=70^{\circ}$， $AC=BC$，以点 $B$ 为旋转中心把 $\triangle ABC$ 按顺时针旋转 $\alpha$ 度，得到 $\triangle A'BC'$，点 $A'$ 恰好落在 $AC$ 上，连接 $CC'$，则 $\angle ACC'=$ \_\_\_\_.\underline{\hspace{2cm}}

\item
  如图， $\triangle ABC$ 的顶点坐标分别为 $A(3, 6)$， $B(1, 3)$， $C(4, 2)$。如果将 $\triangle ABC$ 绕 $C$ 点顺时针旋转 $90^\circ$，得到 $\triangle A'B'C'$，那么点 $A$ 的对应点 $A'$ 的坐标为 \_\_\_\_.\underline{\hspace{2cm}}

  \centerquestionimage{0.35\textwidth}{chuyi_xia_2-p007-cleaned-3.jpg}

\item
  如图，在 Rt$\triangle ABC$中，$\angle$ACB=90$^\circ$，AC=5cm，BC=12cm，将$\triangle ABC$绕点$B$顺时针旋转 60$^\circ$，得到$\triangle BDE$，连接 $DC$交 AB 于点$F$，则$\triangle ACF$与$\triangle BDF$的周长之和为 \_\_\_\_ cm.\underline{\hspace{2cm}}

  \begin{questionimagegroup}
  \questionimageitem{0.32\textwidth}{chuyi_xia_2-p007-cleaned-4.jpg}{}
  \questionimagegap
  \questionimageitem{0.32\textwidth}{chuyi_xia_2-p007-cleaned-5.jpg}{}
  \end{questionimagegroup}

\item
  已知 $\triangle ABC$ 是边长为 $1cm$ 的等边三角形，以 $BC$ 为边作等腰三角形 $BCD$，使得 $DB=DC$，且 $\angle BDC=120^{\circ}$，点 $M$ 是 $AB$ 边上的一个动点，作 $\angle MDN$ 交 $AC$ 边于点 $N$，且满足 $\angle MDN=60^{\circ}$，则 $\triangle AMN$ 的周长为 \_\_\_\_。\underline{\hspace{2cm}}

  \begin{questionimagegroup}
  \questionimageitem{0.32\textwidth}{chuyi_xia_2-p007-cleaned-6.jpg}{}
  \questionimagegap
  \questionimageitem{0.32\textwidth}{chuyi_xia_2-p007-cleaned-7.jpg}{}
  \end{questionimagegroup}

\end{enumerate}

\section*{p008}
\subsection*{四、未分类题目（共 1 小题）}

\solutionlist[14]
\item
  如图，在 $\triangle ABC$ 中， $\angle ACB=90^{\circ}$， $\angle A=20^{\circ}$。在同一平面内，将 $\triangle ABC$ 绕点 $C$ 旋转到 $\triangle A'B'C$ 的位置，设旋转角为 $\alpha$ （$0^{\circ}<\alpha<180^{\circ}$）。若 $\triangle A'B'C$ 中恰有一条边与 $\triangle ABC$ 中的一条边平行，则旋转角 $\alpha$ 的可能的度数为\_\_\_\_。

  \centerquestionimage{0.35\textwidth}{chuyi_xia_2-p008-cleaned-1.jpg}

\end{enumerate}

\subsection*{三、解答题（共 2 小题）}

\solutionlist[15]
\item
  如图，将 $\triangle ABC$ 绕顶点 $A$ 逆时针旋转一角度，使点 $D$ 落在 $BC$ 边上，得到 $\triangle ADE$，此时恰好 $\parallel DE$，已知 $\angle E=35^{\circ}$，求 $\angle DAC$ 的度数.

  \centerquestionimage{0.35\textwidth}{chuyi_xia_2-p008-cleaned-2.jpg}

\item
  如图\circled{1}, 在 $\triangle ABC$ 中, $\angle A=36^{\circ}$ , $AB=AC$ , $\angle ABC$ 的平分线 $BE$ 交 $AC$ 于 $E$ .
(1) 求证: $AE=BC$ ;
（2）如图\circled{2}，过点 $E$ 作 $EF / / BC$ 交 $AB$ 于 $F$，将 $\triangle AEF$ 绕点 $A$ 逆时针旋转角 $\alpha$ （$0^{\circ} < \alpha < 144^{\circ}$）得到 $\triangle AEF$，连接 $CE$， $BF$，求证： $CE=BF$；
（3）在（2）的旋转过程中是否存在 $CE \parallel AB$？若存在，求出相应的旋转角 $\alpha$；若不存在，请说明由.

  \begin{questionimagegroup}
  \questionimageitem{0.28\textwidth}{chuyi_xia_2-p008-cleaned-3.jpg}{图\circled{1}}
  \questionimagegap
  \questionimageitem{0.28\textwidth}{chuyi_xia_2-p008-cleaned-4.jpg}{图\circled{2}}
  \questionimagegap
  \questionimageitem{0.28\textwidth}{chuyi_xia_2-p008-cleaned-5.jpg}{(备用图)}
  \end{questionimagegroup}

\end{enumerate}

\section*{p009}
\subsection*{四、未分类题目（共 1 小题）}

\solutionlist[17]
\item
  两个大小相同且含 $30^{\circ}$ 角的三角板 $ABC$ 和 $DEC$ 如图\circled{1}摆放，使直角顶点重合。将图\circled{1}中 $\triangle DEC$ 绕点 $C$ 逆时针旋转 $30^{\circ}$ 得到图\circled{2}，点 $F$、 $G$ 分别是 $CD$、 $DE$ 与 $AB$ 的交点，点 $H$ 是 $DE$ 与 $AC$ 的交点。
(1) 不添加辅助线, 写出图\circled{2}中所有与 $\triangle BCF$ 全等的三角形;
（2）将图\circled{2}中的$\triangle DEC$绕点$C$逆时针旋转45$^\circ$得$\triangle$ $D_{1}E_{1}C$，点$F$、G、H的对应点分别为 $F_{1}$、 $G_{1}$、 $H_{1}$，如图\circled{3}．探究线段 $D_{1}F_{1}$ 与 $AH_{1}$ 之间的数量关系，并写出推理过程；
(3) 在 (2) 的条件下, 若 $D_{1}E_{1}$ 与 $CE$ 交于点 $I$ , 求证: $G_{1}I=CI$ .

  \begin{questionimagegroup}
  \questionimageitem{0.28\textwidth}{chuyi_xia_2-p009-cleaned-1.jpg}{}
  \questionimagegap
  \questionimageitem{0.28\textwidth}{chuyi_xia_2-p009-cleaned-2.jpg}{图 图\circled{2}}
  \questionimagegap
  \questionimageitem{0.28\textwidth}{chuyi_xia_2-p009-cleaned-3.jpg}{}
  \end{questionimagegroup}

\end{enumerate}

\section*{p010}
\subsection*{一、选择题（共 7 小题）}

\questionlist
\item
  \begin{questionwithimage}{chuyi_xia_2-p010-cleaned-1.jpg}
  下列图形中，既是轴对称图形又是中心对称图形的是（\quad）
  \optionlist
  \item 
  \item 
  \item 
  \item 
  \end{enumerate}
  \end{questionwithimage}

\item
  \begin{questionwithimage}{chuyi_xia_2-p010-cleaned-5.jpg}
  将 $\triangle AOB$ 绕点 $O$ 旋转 $180^{\circ}$ 得到 $\triangle DOE$ , 则下列作图正确的是 （\quad）
  \optionlist
  \item 
  \item 
  \item 
  \item 
  \end{enumerate}
  \end{questionwithimage}

\item
  \begin{questionwithimage}{chuyi_xia_2-p010-cleaned-9.jpg}
  如图， $\triangle ABC$ 与 $\triangle A_1B_1C_1$ 关于点 $O$ 成中心对称，下列说法：\circled{1} $\angle BAC=\angle B_1A_1C_1$；\circled{2} $AC=A_1C_1$；\circled{3} $OA=OA_1$；
\circled{4} $\triangle ABC$ 与 $\triangle A_{1}B_{1}C_{1}$ 的面积相等，其中正确的有（\quad）
  \optionlist
  \item 1个
  \item 2个
  \item 3个
  \item 4个
  \end{enumerate}
  \end{questionwithimage}

\item
  \begin{questionwithimage}{chuyi_xia_2-p010-cleaned-10.jpg}
  如图是一个中心对称图形，A 为对称中心，若 $\angle C=90^{\circ}$， $\angle B=30^{\circ}$， $^{1}BC=1$，则 $BB'$ 的长为（\quad）
  \optionlist
  \item 4
  \item $\frac{\sqrt{3}}{3}$
  \item $\frac{2\sqrt{3}}{3}$
  \item $\frac{4\sqrt{3}}{3}$
  \end{enumerate}
  \end{questionwithimage}

\item
  下列语句中，不正确的是（\quad）
  \optionlist
  \item 图形平移是由移动的方向和距离所决定的
  \item 图形旋转是由旋转中心和旋转角度所决定的
  \item 任意两点都成中心对称
  \item 任意两条相等的线段都成中心对称
  \end{enumerate}

\item
  下列描述中心对称的特征的语句中，其中正确的是（\quad）
  \optionlist
  \item 成中心对称的两个图形中，连接对称点的线段不一定经过对称中心
  \item 成中心对称的两个图形中，对称中心不一定平分连接对称点的线段
  \item 成中心对称的两个图形中，对称点的连线一定经过对称中心，但不一定被对称中心平分
  \item 成中心对称的两个图形中，对称点的连线一定经过对称中心，且被对称中心平分
  \end{enumerate}

\item
  一个正多边形绕它的中心旋转 $45^{\circ}$ 后，就与原正多边形第一次重合，那么这个正多边形（\quad）
  \optionlist
  \item 是轴对称图形，但不是中心对称图形
  \item 是中心对称图形，但不是轴对称图形
  \item 既是轴对称图形，又是中心对称图形
  \item 既不是轴对称图形，也不是中心对称图形
  \end{enumerate}

\end{enumerate}

\section*{p011}
\subsection*{二、填空题（共 7 小题）}

\blanklist[8]
\item
  如图，是一个中心对称图形，A 为对称中心，若 $\angle C=90^{\circ}$， $\angle B=30^{\circ}$，AC=1，则 $BB'$ 的长为 \_\_\_\_.\underline{\hspace{2cm}}

  \centerquestionimage{0.35\textwidth}{chuyi_xia_2-p011-cleaned-1.jpg}

\item
  $\triangle ABC$ 与 $\triangle A'B'C'$ 关于原点 $O$ 成中心对称，点 $A$、 $B$、 $C$ 的对称点分别是 $A'$、 $B'$、 $C'$，若 $AB=3$， $AC=1$，则 $B'C'$ 的范围是 \_\_\_\_.\underline{\hspace{2cm}}

\item
  如图，是 $4\times 4$ 的正方形网格，把其中一个标有数字的白色小正方形涂黑就可以使图中的黑色部分构成一个中心对称图形，则这个白色小正方形内的数字是\_\_\_\_。\underline{\hspace{2cm}}

  \centerquestionimage{0.35\textwidth}{chuyi_xia_2-p011-cleaned-2.jpg}

\item
  在平面直角坐标系中，点 $P(1,1)$， $N(2,0)$， $\triangle MNP$ 和 $\triangle M_1N_1P_1$ 的顶点都在格点上， $\triangle MNP$ 与 $\triangle M_1N_1P_1$ 是关于某一点中心对称，则对称中心的坐标为 \_\_\_\_.\underline{\hspace{2cm}}

  \centerquestionimage{0.35\textwidth}{chuyi_xia_2-p011-cleaned-3.jpg}

\item
  \begin{questionwithimage}{chuyi_xia_2-p011-cleaned-4.jpg}
  将五个边长都为 3cm 的正方形按如图所示摆放，点 $A$、$B$、$C$、$D$ 分别是四个正方形的中心，则图中四块阴影部分的面积和是 \_\_\_\_ cm$^2$。
  \end{questionwithimage}

\item
  在平面直角坐标系中，点 $A, B, C$ 的坐标分别为 $(1, 0),(0, 1),(-1, 0)$。一个电动玩具从坐标原点 $O$ 出发，第一次跳跃到点 $P_{1}$。使得点 $P_{1}$ 与点 $O$ 关于点 $A$ 成中心对称；第二次跳跃到点 $P_{2}$，使得点 $P_{2}$ 与点 $P_{1}$ 关于点 $B$ 成中心对称；第三次跳跃到点 $P_{3}$，使得点 $P_{3}$ 与点 $P_{2}$ 关于点 $C$ 成中心对称；第四次跳跃到点 $P_{4}$，使得点 $P_{4}$ 与点 $P_{3}$ 关于点 $A$ 成中心对称；第五次跳跃到点 $P_{5}$，使得点 $P_{5}$ 与点 $P_{4}$ 关于点 $B$ 成中心对称；…照此规律重复下去，则点 $P_{2019}$ 的坐标为 \_\_\_\_。\underline{\hspace{2cm}}

\item
  如果将点 $P$ 绕定点 $M$ 旋转 $180^{\circ}$ 后与点 $Q$ 重合，那么称点 $P$ 与点 $Q$ 关于点 $M$ 对称，定点 $M$ 叫做对称中心。此时，点 $M$ 是线段 $PQ$ 的中点。在平面直角坐标系中， $\triangle ABO$ 的顶点 $A, B, O$ 的坐标分别为（1，0）、（0，1）、（0，0）。点列 $P_{1}, P_{2}, P_{3}, \ldots$，中的相邻两点都关于 $\triangle ABO$ 的一个顶点对称：点 $P_{1}$ 与点 $P_{2}$ 关于点 $A$ 对称，点 $P_{2}$ 与点 $P_{3}$ 关于点 $B$ 对称，点 $P_{3}$ 与点 $P_{4}$ 关于点 $O$ 对称，点 $P_{4}$ 与点 $P_{5}$ 关于点 $A$ 对称，点 $P_{5}$ 与点 $P_{6}$ 关于点 $B$ 对称，点 $P_{6}$ 与点 $P_{7}$ 关于点 $O$ 对称， $\ldots$，对称中心分别是 $A, B, O, A, B, O, \ldots$，且这些对称中心依次循环。已知点 $P_{1}$ 的坐标是（1，1），则点 $P_{2019}$ 的坐标为 \_\_\_\_。\underline{\hspace{2cm}}

  \centerquestionimage{0.35\textwidth}{chuyi_xia_2-p011-cleaned-5.jpg}

\end{enumerate}

\section*{p012}
\subsection*{三、解答题（共 4 小题）}

\solutionlist[15]
\item
  在 $14\times 9$ 的方格纸中, 每个小正方形的边长都为 1, $\triangle ABC$ 与 $\triangle A'B'C'$ 的位置如图所示;
（1）请说明$\triangle ABC$与$\triangle$A'B'C'的位置关系；
（2）若点$C$的坐标为(0，0)，则点 $B'$ 的坐标为\_\_\_\_；
(3) 求线段$CC$的长.

  \centerquestionimage{0.35\textwidth}{chuyi_xia_2-p012-cleaned-1.jpg}

\item
  已知 $MN \perp PQ$ 于点 $O$，点 $A_{1}$ 和点 $A$ 关于 $MN$ 对称，点 $A_{2}$ 和点 $A$ 关于 $PQ$ 对称，试证明：点 $A_{1}$ 和点 $A_{2}$ 关于点 $O$ 成中心对称.
(1) 求证: AC = CD;

\item
  已知：如图，三角形 $ABM$ 与三角形 $ACM$ 关于直线 $AF$ 成轴对称，三角形 $ABE$ 与三角形 $DCE$ 于点 $E$ 成中心对称，点 $E$、 $D$、 $M$ 都在线段 $AF$ 上， $BM$ 的延长线交 $CF$ 于点 $P$ .
(2) 若 $\angle BAC=2\angle MPC$，请你判断 $\angle F$ 与 $\angle MCD$ 的数量关系，并说明理由.

  \centerquestionimage{0.35\textwidth}{chuyi_xia_2-p012-cleaned-2.jpg}

\item
  \begin{questionwithimage}{chuyi_xia_2-p012-cleaned-3.jpg}
  课外兴趣小组活动时，老师提出了如下问题：
（1）如图\circled{1}，在$\triangle ABC$中，若AB=5，AC=3，求$BC$边上的中线$AD$的取值范围。小明在组内经过合作交流，得到了如下的解决方法：延长 $AD$到E，使得DE=AD，将$\triangle ACD$绕点$D$逆时针旋转180$^\circ$得到$\triangle EBD$)，把A（\quad）
  \optionlist
  \item A
  \item 2AD集中在$\triangle ABE$中，利用关系可得2<AE<8，则1<AD<4。
[感悟]解题时，条件中若出现“中点”“中线”字样，可以考虑构造以中点为对称中心的把分散的已知条件和所求证的结论集中到同一个三角形中.
（2）解决问题：受到（1）的启发，请你证明下列命题：如图\circled{2}，在$\triangle ABC$中，D是 $DE \perp DF$，DE交AB于点$E$，DF交AC于点$F$，连接 $EF$.
求证： $BE+CF>EF$，若 $\angle A=90^{\circ}$，探索线段$BE$、C
  \item EF之间的等量关系，并加
  \end{enumerate}

  % 图注：图\circled{1}
  \end{questionwithimage}

\end{enumerate}

\section*{p013}
\subsection*{四、未分类题目（共 1 小题）}

\solutionlist[18]
\item
  课外兴趣小组活动时, 老师提出了如下问题:
(1) 如图\circled{1}, 在 $\triangle ABC$ 中, 若 $AB=5$ , $AC=3$ , 求 $BC$ 边上的中线 $AD$ 的取值范围.
小明在组内经过合作交流，得到了如下的解决方法：延长 $AD$ 到 $E$，使得 $DE=AD$，再连接 $BE$ （或将 $\triangle ACD$ 绕点 $D$ 逆时针旋转 $180^{\circ}$ 得到 $\triangle EBD$ )，把 $AB$、 $AC$、 $2AD$ 集中在 $\triangle ABE$ 中，利用三角形的三边关系可得 $2 < AE < 8$，则 $1 < AD < 4$
[感悟]解题时，条件中若出现“中点”“中线”字样，可以考虑构造以中点为对称中心的中心对称图形，把分散的已知条件和所求证的结论集中到同一个三角形中.
（2）解决问题：受到（1）的启发，请你证明下列命题：如图\circled{2}，在 $\triangle ABC$ 中， $D$ 是 $BC$ 边上的中点， $DE \perp DF$， $DE$ 交 $AB$ 于点 $E$， $DF$ 交 $AC$ 于点 $F$，连接 $EF$
求证： $BE+CF>EF$，若 $\angle A=90^{\circ}$，探索线段$BE$、CF、EF 之间的等量关系，并加以证明.

  \begin{questionimagegroup}
  \questionimageitem{0.32\textwidth}{chuyi_xia_2-p013-cleaned-1.jpg}{图\circled{1}}
  \questionimagegap
  \questionimageitem{0.32\textwidth}{chuyi_xia_2-p013-cleaned-2.jpg}{图\circled{2}}
  \end{questionimagegroup}

\end{enumerate}

\section*{p014}
\subsection*{一、选择题（共 7 小题）}

\questionlist
\item
  点 $A(3,-1)$ 关于原点的对称点 $A'$ 的坐标是（\quad）
  \optionlist
  \item $(-3,-1)$
  \item $(3, 1)$
  \item $(-3, 1)$
  \item $(-1, 3)$
  \end{enumerate}

\item
  在直角坐标系中，将点（-2，3）关于原点的对称点向左平移2个单位长度得到的点的坐标是（\quad）
  \optionlist
  \item (4, -3)
  \item (-4, 3)
  \item (0, -3)
  \item (0, 3)
  \end{enumerate}

\item
  在平面直角坐标系中，对于平面内任一点 $(a, b)$，若规定以下三种变换： $\circled{1} \triangle(a, b)=(-a, b)$； $\circled{2} \mathsf{O}(a, b)=(-a,-b)$； $\circled{3} \Omega(a, b)=(a,-b)$；
按照以上变换有： $\triangle(O(1,2))=(1,-2)$，那么 $O(\Omega(3,4))$ 等于（\quad）
  \optionlist
  \item $(3,4)$
  \item $(3,-4)$
  \item $(-3,4)$
  \item $(-3,-4)$
  \end{enumerate}

\item
  在直角坐标系中，点 $A$ 的坐标为 $(-3, 4)$，那么下列说法正确的是（\quad）
  \optionlist
  \item 点 $A$ 与点 $B$ （-3，-4）关于 $y$ 轴对称
  \item 点 $A$ 与点 $C$ （3，-4）关于 $x$ 轴对称
  \item 点 $A$ 与点 $C$ （4，-3）关于原点对称
  \item 点 $A$ 与点 $F$ （-4，3）关于第二象限的平分线对称
  \end{enumerate}

\item
  在平面直角坐标系中，关于点 $A(\sqrt{3},-1)$ 的图象变化有以下说法： $\circled{1}$ 点 $A$ 关于 $y$ 轴的对称点 $B$ 的坐标为 $(-\sqrt{3},-1)$ $\circled{2}$ 点 $A$ 与点 $C(-1, \sqrt{3})$ 关于原点对称 $\circled{3}$ 把点 $A$ 先向右平移2个单位长度，再向下平移3个单位长度得到点 $D(2+\sqrt{3},-4)$ $\circled{4}$ 把点 $A$ 绕原点顺时针旋转 $30^{\circ}$，得到点 $E(1,-\sqrt{3})$ 其中，正确的说法是（\quad）
  \optionlist
  \item $\circled{1}\circled{3}\circled{4}$
  \item $\circled{1}\circled{2}\circled{3}\circled{4}$
  \item $\circled{1}\circled{2}\circled{3}$
  \item $\circled{2}\circled{3}\circled{4}$
  \end{enumerate}

\item
  \begin{questionwithimage}{chuyi_xia_2-p014-cleaned-1.jpg}
  如图，把图中的 $\triangle ABC$ 经过一定的变换得到 $\triangle A'B'C'$，如果图中 $\triangle ABC$ 上的点 $P$ 的坐标为 $(a, b)$，那么它的对应点 $P'$ 的坐标为（\quad）
  \optionlist
  \item $(a-2, b)$
  \item $(a+2, b)$
  \item $(-a-2,-b)$
  \item $(a+2,-b)$
  \end{enumerate}
  \end{questionwithimage}

\item
  \begin{questionwithimage}{chuyi_xia_2-p014-cleaned-2.jpg}
  如图，平面直角坐标系， $\angle ABO=90^\circ$，将直角 $\triangle AOB$ 绕 $O$ 点顺时针旋转，使点 $B$ 落在 $x$ 轴上的点 $B_1$ 处，点 $A$ 落在 $A_1$ 处，若 $B$ 点的坐标为 $\left(\frac{16}{5}, \frac{12}{5}\right)$，则线段 $AA_1$ 的长度是（\quad）
  \optionlist
  \item $2\sqrt{2}$
  \item $\sqrt{10}$
  \item $2\sqrt{7}$
  \item $2\sqrt{5}$
  \end{enumerate}
  \end{questionwithimage}

\end{enumerate}

\section*{p015}
\subsection*{二、填空题（共 7 小题）}

\blanklist[9]
\item
  如图，Rt$\triangle OAB$的直角边$OA$在 y 轴上，点$B$在第一象限内，OA=2，AB=1，若将$\triangle OAB$绕点$O$按逆时针方向旋转 90$^\circ$，则点$B$的对应点的坐标为 \_\_\_\_.\underline{\hspace{2cm}}

  \centerquestionimage{0.35\textwidth}{chuyi_xia_2-p015-cleaned-1.jpg}

\item
  在平面直角坐标系中，已知点 $O(0,0), A(1,n), B(2,0)$，其中 $n > 0$， $\triangle OAB$ 是等边三角形。点 $P$ 是线段 $OB$ 的中点，将 $\triangle OAB$ 绕点 $O$ 逆时针旋转 $30^\circ$，记点 $P$ 的对应点为点 $Q$，则 $n=$ \_\_\_\_，点 $Q$ 的坐标是 \_\_\_\_。\underline{\hspace{2cm}}

\item
  根据指令 $[s, A]$ ($s \geqslant 0, 0^{\circ} < A < 180^{\circ}$ ), 机器人在平面上能完成下列动作: 先原地逆时针旋转角度 $A$ , 再朝其面对的方向沿直线行走距离 $s$ , 现机器人在直角坐标系的坐标原点, 且面对 $x$ 轴正方向.
(1) 若给机器人下了一个指令[4, 60$^\circ$], 则机器人应移动到点\_\_\_\_;
(2) 请你给机器人下一个指令\_\_\_\_，使其移动到点 $(-5, 5)$ .\underline{\hspace{2cm}}

\item
  在一个坐标方格盘中，你可操纵一只遥控机器蛙在方格盘上进行跳步游戏，机器蛙每次跳步只能按如下两种方式（第一种：向上、下、左、右可任意跳动1格或3格；第二种跳到关于原点的对称点上）中的一种进行。若机器蛙在点 $A(-5, 4)$，现欲操纵它跳到点 $B(2,-3)$，请问机器蛙至少要跳一次。\underline{\hspace{2cm}}

\item
  如图，在直角坐标系中，四边形 $ABCD$ 是正方形，点 $A$ 的坐标为（0，1），点 $B$ 的坐标为（2，0）.
\circled{1}点$C$ 的坐标为\_\_\_\_；
\circled{2}若正方形 ABCD 和正方形 $A_{1}BC_{1}B_{1}$ 关于点$B$成中心对称；正方形 $A_{1}BC_{1}B_{1}$ 和正方形 $A_{2}B_{2}C_{2}B_{1}$ 关于点 $B_{1}$ 成中心对称；…，依此规律，则点 $C_{6}$ 的坐标为 \_\_\_\_。\underline{\hspace{2cm}}

  \centerquestionimage{0.35\textwidth}{chuyi_xia_2-p015-cleaned-2.jpg}

\item
  已知，正六边形 $ABCDEF$ 在直角坐标系内的位置如图所示， $A(-2,0)$，点 $B$ 在原点，把正六边形 $ABCDEF$ 沿 $x$ 轴正半轴作无滑动的连续翻转，每次翻转 $60^\circ$，经过 2019 次翻转之后，点 $B$ 的坐标是\_\_\_\_.
·三. 解答题（共5小题）\underline{\hspace{2cm}}

  \centerquestionimage{0.35\textwidth}{chuyi_xia_2-p015-cleaned-3.jpg}

\item
  \begin{questionwithimage}{chuyi_xia_2-p015-cleaned-4.jpg}
  每个小方格都是边长为 1 个单位长度的正方形，在建立平面直角坐标系后， $\triangle ABC$ 的顶点均在格点上，
\circled{1}写出\underline{\hspace{2cm}}（\quad）
  \optionlist
  \item 
  \item C 的坐标.
\circled{2}以原点 $O$ 为对称中心，画出 $\triangle ABC$ 关于原点 $O$ 对称的 $\triangle A_{1}B_{1}C_{1}$，并写出 $A_{1}, B_{1}, C_{1}$ .
  \end{enumerate}
  \end{questionwithimage}

\end{enumerate}

\section*{p016}
\subsection*{四、未分类题目（共 3 小题）}

\solutionlist[16]
\item
  已知：如图， $\triangle ABC$、直线 $m$、点 $M$ 在网格中如图所示的位置，请按以下要求作图：
(1) 将 $\triangle ABC$ 向上平移 6 个单位得 $\triangle A_{1}B_{1}C_{1}$ ;
(2) 作出 $\triangle ABC$ 关于直线 $m$ 的轴对称图形 ${\Delta }_{{A}_{2}{B}_{2}{C}_{2}}$ ;
(3) 作出 $\triangle A_{2}B_{2}C_{2}$ 绕点$M$顺时针旋转 $90^{\circ}$ 的图形 $\triangle A_{3}B_{3}C_{3}$ .

  \centerquestionimage{0.35\textwidth}{chuyi_xia_2-p016-cleaned-1.jpg}

\item
  在平面直角坐标系中，对于平面内任一点 $(m, n)$，规定以下三种变换：
\circled{1} $f(m, n)=(m,-n)$ ; \circled{2} $g(m, n)=(-m, n)$ ; \circled{3} $h(m, n)=(-m,-n)$ .
（1）请你根据以上规定的变换，求 $f[g(-3,2)]$ 的值；
(2) 请你以点 $(a, b)$ 为例, 探索以上三种变换之间的关系.

\item
  如图，在平面直角坐标系中，已知 $\triangle AOB$ 是等边三角形，点 $A$ 的坐标是（0，3），点 $B$ 在第一象限， $\angle OAB$ 的平分线交 $x$ 轴于点 $P$，把 $\triangle AOP$ 绕着点 $A$ 按逆时针方向旋转，使边 $AO$ 与 $AB$ 重合，得 $\triangle ABD$，连接 $DP$。求： $DP$ 的长及点 $D$ 的坐标。

  \centerquestionimage{0.35\textwidth}{chuyi_xia_2-p016-cleaned-2.jpg}

\end{enumerate}

\section*{p017}
\subsection*{四、未分类题目（共 1 小题）}

\solutionlist[19]
\item
  如图，正方形 $ABCD$ 与正方形 $A_{1}B_{1}C_{1}D_{1}$ 关于某点中心对称，已知 $A$、$D_{1}$、$D$ 三点的坐标分别是 $(0, 4)$、$(0, 3)$、$(0, 2)$。
  (1) 求对称中心的坐标。
  (2) 写出顶点 $B$、$C$、$B_{1}$、$C_{1}$ 的坐标。

  \centerquestionimage{0.35\textwidth}{chuyi_xia_2-p017-cleaned-1.jpg}

\end{enumerate}

\section*{p018}
\subsection*{四、未分类题目（共 3 小题）}

\solutionlist[3]
\item
  如图，在四边形 ABCD 中， $\angle ABC=\angle ADC=120^{\circ}$，AB = BC = CD = DA，点$E$是点$D$，直线$BD$与直线$AC$交于点$E$，连接 $EH$.
点，连接 $DE$， $\angle DEC=50^{\circ}$，将线段$BC$绕点$B$逆时针旋转 $50^{\circ}$ 并延长得到射线$BF$，交点$G$.
(1) 如图\circled{1}, 当 $\angle BAC$ 为锐角时,
\circled{1}求证： $BE \perp AC;$
\circled{2}求$\angle$BEH的度数；
（2）当 $\angle BAC$ 为钝角时，请依题意用实线补全图\circled{2}，并直接写出表示线段 $EC$，ED，EH之间的数量关系的等式.
(1) 依题意补全图形;

  \begin{questionimagegroup}
  \questionimageitem{0.32\textwidth}{chuyi_xia_2-p018-cleaned-1.jpg}{图\circled{1}}
  \questionimagegap
  \questionimageitem{0.32\textwidth}{chuyi_xia_2-p018-cleaned-2.jpg}{图\circled{2}}
  \end{questionimagegroup}

\item
  在 $\triangle ABC$ 中， $\angle C=90^{\circ}$， $AC=BC$，点 $D$ 在射线 $BC$ 上（不与点 $B$、 $C$ 重合），连接 $AD$，将 $AD$ 绕点时针旋转 $90^{\circ}$ 得到 $DE$，连接 $BE$ .
(1) 如图\circled{1}, 点$D$在 $BC$边上.
\circled{1}依题意补全图\circled{1};
\circled{2}作 $DF \perp BC$ 交 $AB$ 于点 $F$，若 $AC=8$， $DF=3$，求 $BE$ 的长；
（2）如图\circled{2}，点 $D$ 在 $BC$ 边的延长线上，用等式表示线段 $AB$、 $BD$、 $BE$ 之间的数量关系（直接写结论）.
(2) 求证: $EG=BC$ ;
(3) 用等式表示线段$AE$，EG，BG 之间的数量关系：\_\_\_\_

  \begin{questionimagegroup}
  \questionimageitem{0.28\textwidth}{chuyi_xia_2-p018-cleaned-3.jpg}{图\circled{1}}
  \questionimagegap
  \questionimageitem{0.28\textwidth}{chuyi_xia_2-p018-cleaned-4.jpg}{}
  \questionimagegap
  \questionimageitem{0.28\textwidth}{chuyi_xia_2-p018-cleaned-5.jpg}{}
  \questionimagegap
  \questionimageitem{0.28\textwidth}{chuyi_xia_2-p018-cleaned-6.jpg}{名}
  \end{questionimagegroup}

\item
  在 Rt$\triangle ABC$中，$\angle$ACB=90$^\circ$，D 是 AB 的中点，DE$\perp$BC 于 E，连接 $CD$
(1) 如图\circled{1}, 如果 $\angle A=30^{\circ}$ , 那么 $DE$ 与 $CE$ 之间的数量关系是
（2）如图\circled{2}，在（1）的条件下，P是线段$CB$上一点，连接 $DP$，将线段$DP$得到线段$DF$，连接 $BF$，请猜想DE、BF、BP三者之间的数量关系，并证明你在

  \begin{questionimagegroup}
  \questionimageitem{0.35\textwidth}{chuyi_xia_2-p018-cleaned-7.jpg}{图\circled{1}}
  \end{questionimagegroup}

\end{enumerate}

\section*{p019}
\subsection*{四、未分类题目（共 2 小题）}

\solutionlist[3]
\item
  如图，在四边形 $ABCD$ 中，$\angle ABC=\angle ADC=120^\circ$，$AB=BC=CD=DA$，点 $E$ 是对角线 $AC$ 上一点，连接 $DE$，$\angle DEC=50^\circ$。将线段 $BC$ 绕点 $B$ 逆时针旋转 $50^\circ$ 并延长得到射线 $BF$，交 $ED$ 的延长线于点 $G$。
  (1) 依题意补全图形；
  (2) 求证：$EG=BC$；
  (3) 用等式表示线段 $AE$、$EG$、$BG$ 之间的数量关系：\_\_\_\_。

  \begin{questionimagegroup}
  \questionimageitem{0.32\textwidth}{chuyi_xia_2-p019-cleaned-1.jpg}{}
  \questionimagegap
  \questionimageitem{0.32\textwidth}{chuyi_xia_2-p019-cleaned-2.jpg}{备用图}
  \end{questionimagegroup}

\item
  在 $Rt\triangle ABC$ 中，$\angle ACB=90^\circ$，$D$ 是 $AB$ 的中点，$DE \perp BC$ 于点 $E$，连接 $CD$。
  (1) 如图\circled{1}，如果 $\angle A=30^\circ$，那么 $DE$ 与 $CE$ 之间的数量关系是 \_\_\_\_；
  (2) 如图\circled{2}，在（1）的条件下，$P$ 是线段 $CB$ 上一点，连接 $DP$，将线段 $DP$ 绕点 $D$ 逆时针旋转 $60^\circ$ 得到线段 $DF$，连接 $BF$，请猜想 $DE$、$BF$、$BP$ 三者之间的数量关系，并证明你的结论。

  \begin{questionimagegroup}
  \questionimageitem{0.32\textwidth}{chuyi_xia_2-p019-cleaned-3.jpg}{图\circled{1}}
  \questionimagegap
  \questionimageitem{0.32\textwidth}{chuyi_xia_2-p019-cleaned-4.jpg}{图\circled{2}}
  \end{questionimagegroup}

\end{enumerate}

\section*{p020}
\subsection*{四、未分类题目（共 1 小题）}

\solutionlist[5]
\item
  如图\circled{1}，已知线段 $BC=2$，点 $B$ 关于直线 $AC$ 的对称点是点 $D$，点 $E$ 为射线 $CA$ 上一点，连接 $DE$、$BE$。
  (1) 依题意补全图\circled{1}，并证明：$\triangle BDE$ 为等边三角形；
  (2) 若 $\angle ACB=45^\circ$，点 $C$ 关于直线 $BD$ 的对称点为点 $F$，连接 $FD$、$FB$。将 $\triangle CDE$ 绕点 $D$ 旋转 $\alpha$ 度（$0^\circ<\alpha<360^\circ$）得到 $\triangle C'DE'$，点 $E$ 的对应点为 $E'$，点 $C$ 的对应点为点 $C'$。
  \circled{1} 如图\circled{2}，当 $\alpha=30^\circ$ 时，连接 $BC'$。证明：$EF=BC'$；
  \circled{2} 如图\circled{3}，点 $M$ 为 $DC$ 中点，点 $P$ 为线段 $C'E'$ 上的任意一点，试探究在此旋转过程中相关线段长度的取值范围。

  \begin{questionimagegroup}
  \questionimageitem{0.28\textwidth}{chuyi_xia_2-p020-cleaned-1.jpg}{图\circled{1}}
  \questionimagegap
  \questionimageitem{0.28\textwidth}{chuyi_xia_2-p020-cleaned-2.jpg}{图\circled{2}}
  \questionimagegap
  \questionimageitem{0.28\textwidth}{chuyi_xia_2-p020-cleaned-3.jpg}{图\circled{3}}
  \end{questionimagegroup}

\end{enumerate}

\section*{p021}
\subsection*{四、未分类题目（共 2 小题）}

\solutionlist
\item
  如图\circled{1}，在 $Rt\triangle ABC$ 中，$\angle ACB=90^\circ$，$E$ 是边 $AC$ 上任意一点（点 $E$ 与点 $A$、$C$ 不重合），以 $CE$ 为一直角边作 $Rt\triangle ECD$，且 $\angle ECD=90^\circ$，连接 $BE$、$AD$。
  (1) 若 $CA=CB$，$CE=CD$：
  \circled{1} 猜想线段 $BE$、$AD$ 之间的数量关系及所在直线的位置关系，直接写出结论；
  \circled{2} 现将图\circled{1}中的 $Rt\triangle ECD$ 绕着点 $C$ 顺时针旋转锐角 $\alpha$，得到图\circled{2}，请判断 \circled{1} 中的结论是否依然成立；若成立，请证明；若不成立，请说明理由。
  (2) 若图\circled{2}中 $CA=CB=7$，$CE=CD=4$，连接 $BD$、$AE$，计算 $BD^2+AE^2$ 的值。

  \begin{questionimagegroup}
  \questionimageitem{0.32\textwidth}{chuyi_xia_2-p021-cleaned-1.jpg}{图\circled{1}}
  \questionimagegap
  \questionimageitem{0.32\textwidth}{chuyi_xia_2-p021-cleaned-2.jpg}{图\circled{2}}
  \end{questionimagegroup}

\item
  如图，$\triangle ABC$ 中，$\angle BAC=90^\circ$，$AB=AC$，边 $BA$ 绕点 $B$ 顺时针旋转 $\alpha$ 角得到线段 $BP$，连接 $PA$、$PC$，过点 $P$ 作 $PD \perp AC$ 于点 $D$。
  (1) 如图\circled{1}，若 $\alpha=60^\circ$，求 $\angle DPC$ 的度数；
  (2) 如图\circled{2}，若 $\alpha=30^\circ$，直接写出 $\angle DPC$ 的度数；
  (3) 如图\circled{3}，若 $\alpha=150^\circ$，依题意补全图，并求 $\angle DPC$ 的度数。

  \begin{questionimagegroup}
  \questionimageitem{0.28\textwidth}{chuyi_xia_2-p021-cleaned-3.jpg}{图\circled{1}}
  \questionimagegap
  \questionimageitem{0.28\textwidth}{chuyi_xia_2-p021-cleaned-4.jpg}{图\circled{2}}
  \questionimagegap
  \questionimageitem{0.28\textwidth}{chuyi_xia_2-p021-cleaned-5.jpg}{图\circled{3}}
  \end{questionimagegroup}

\end{enumerate}

\section*{p022}
\subsection*{四、未分类题目（共 2 小题）}

\solutionlist[3]
\item
  如图\circled{1}，在 $\triangle ABC$ 中，$CA=CB$，$\angle ACB=90^\circ$，$D$ 是 $\triangle ABC$ 内部一点，$\angle ADC=135^\circ$。将线段 $CD$ 绕点 $C$ 逆时针旋转 $90^\circ$ 得到线段 $CE$，连接 $DE$。
  (1) \circled{1} 依题意补全图形；
  \circled{2} 请判断 $\angle ADC$ 和 $\angle CDE$ 之间的数量关系，并直接写出答案。
  (2) 在（1）的条件下，连接 $BE$，过点 $C$ 作 $CM \perp DE$，请判断线段 $CM$、$AE$ 和 $BE$ 之间的数量关系并说明理由。
  (3) 如图\circled{2}，在正方形 $ABCD$ 中，$AB=\sqrt{2}$，如果 $PD=1$，$\angle BPD=90^\circ$，请直接写出点 $A$ 到 $BP_{1}$ 的距离。

  \begin{questionimagegroup}
  \questionimageitem{0.32\textwidth}{chuyi_xia_2-p022-cleaned-1.jpg}{图\circled{1}}
  \questionimagegap
  \questionimageitem{0.32\textwidth}{chuyi_xia_2-p022-cleaned-2.jpg}{图\circled{2}}
  \end{questionimagegroup}

\item
  已知 $\triangle ABC$ 是锐角三角形，$BA=BC$，点 $E$ 为 $AC$ 边的中点，点 $D$ 为 $AB$ 边上一点，且 $\angle ABC=\angle AED=\alpha$。
  (1) 如图\circled{1}，当 $\alpha=40^\circ$ 时，$\angle ADE=$ \_\_\_\_ $^\circ$；
  (2) 如图\circled{2}，取 $BC$ 边的中点 $F$，连接 $FD$，将 $\angle AED$ 绕点 $E$ 顺时针旋转适当的角度 $\beta$（$\beta<\alpha$），得到相应图形，$EM$ 与 $BA$ 的延长线交于点 $M$，$EN$ 与 $FD$ 的延长线交于点 $N$。
  \circled{1} 依题意补全图形；
  \circled{2} 猜想线段 $EM$ 与 $EN$ 之间的数量关系，并证明你的结论。

  \begin{questionimagegroup}
  \questionimageitem{0.28\textwidth}{chuyi_xia_2-p022-cleaned-3.jpg}{图\circled{1}}
  \questionimagegap
  \questionimageitem{0.28\textwidth}{chuyi_xia_2-p022-cleaned-4.jpg}{图\circled{2}}
  \questionimagegap
  \questionimageitem{0.28\textwidth}{chuyi_xia_2-p022-cleaned-5.jpg}{}
  \end{questionimagegroup}

\end{enumerate}

\section*{p023}
\subsection*{四、未分类题目（共 2 小题）}

\solutionlist[5]
\item
  在 $\triangle ABC$ 中，$AB=BC=2$，$\angle ABC=90^\circ$，$BD$ 为斜边 $AC$ 上的中线。将 $\triangle ABD$ 绕点 $D$ 顺时针旋转 $\alpha$（$0^\circ<\alpha<180^\circ$）得到 $\triangle EFD$，其中点 $A$ 的对应点为点 $E$，点 $B$ 的对应点为点 $F$，$BE$ 与 $FC$ 相交于点 $H$。
  (1) 如图\circled{1}，直接写出 $BE$ 与 $FC$ 的数量关系；
  (2) 如图\circled{2}，$M$、$N$ 分别为 $EF$、$BC$ 的中点。求证：$MN=\frac{\sqrt{2}}{2}FC$；
  (3) 连接 $BF$、$CE$，如图\circled{3}，直接写出在此旋转过程中，线段 $BF$、$CE$ 与 $AC$ 之间的数量关系：\_\_\_\_。

  \begin{questionimagegroup}
  \questionimageitem{0.28\textwidth}{chuyi_xia_2-p023-cleaned-1.jpg}{图\circled{1}}
  \questionimagegap
  \questionimageitem{0.28\textwidth}{chuyi_xia_2-p023-cleaned-2.jpg}{图\circled{2}}
  \questionimagegap
  \questionimageitem{0.28\textwidth}{chuyi_xia_2-p023-cleaned-3.jpg}{图\circled{3}}
  \end{questionimagegroup}

\item
  \begin{questionwithimage}{chuyi_xia_2-p023-cleaned-4.jpg}
  数学活动课上，老师提出这样一个问题: 如果 AB=BC， $\angle ABC=60^{\circ}$， $\angle APC=30^{\circ}$，连接 $PB$，那么 P（\quad）
  \optionlist
  \item P
  \item PC 之间会有怎样的等量关系呢？经过思考后，部分同学进行了如下的交流：
小蕾：我将图形进行了特殊化，让点 $P$ 在 $BA$ 延长线上（如图\circled{1}），得到了一个猜想： $PA^{2}+PC^{2}=PB^{2}$
小东：我假设点 $P$ 在 $\angle ABC$ 的内部，根据题目条件，这个图形具有“共端点等线段”的特点，可以利用旋转解决问题，旋转 $\triangle PAB$ 后得到 $\triangle P'CB,$ 并且可推出 $\triangle PBP',\triangle PCP'$ 分别是等边三角形、直角三角形，就能得到猜想和证明方法.
(1) 如图\circled{2}, 点$P$在 $\angle ABC$ 的内部,
\circled{1} ${PA}=4,{PC}=2\sqrt{3},{PB}=$ \_\_\_\_. \circled{2}用等式表示 ${PA}$、 ${PB}$、 ${PC}$ 之间的数量关系,并证明.
(2) 对于点$P$的其他位置，是否始终具有\circled{2}中的结论？若是，请证明；若不是，请举例说明.
  \end{enumerate}
  \end{questionwithimage}

\end{enumerate}

\section*{p024}
\subsection*{一、选择题（共 6 小题）}

\questionlist
\item
  在平移、旋转和轴对称这些图形变换下，它们共同具有的特征是（\quad）
  \optionlist
  \item 图形的形状、大小没有改变，对应线段平行且相等
  \item 图形的形状、大小没有改变，对应线段垂直，对应角相等
  \item 图形的形状、大小都发生了改变，对应线段相等，对应角相等
  \item 图形的形状、大小没有改变，对应线段相等，对应角相等
  \end{enumerate}

\item
  \begin{questionwithimage}{chuyi_xia_2-p024-cleaned-1.jpg}
  在方格纸中，图\circled{1}中的图形 $N$ 经过旋转平移后的位置如图\circled{2}所示，那么下列说法正确的是（\quad）
  \optionlist
  \item 绕 $A$ 点顺时针旋转 $90^{\circ}$，再向下平移3个单位
  \item 绕 $A$ 点逆时针旋转 $90^{\circ}$，再向下平移3个单位
  \item 绕 $A$ 点顺时针旋转 $90^{\circ}$，再向下平移5个单位
  \item 绕 $A$ 点逆时针旋转 $90^{\circ}$，再向下平移4个单位
  \end{enumerate}

  % 图注：图\circled{1}
  \end{questionwithimage}

\item
  \begin{questionwithimage}{chuyi_xia_2-p024-cleaned-3.jpg}
  在所示的四个三角形中, 不能与 $\triangle ABC$ 经过旋转变换得到的是 （\quad）
  \optionlist
  \item 
  \item 
  \item 
  \end{enumerate}
  \end{questionwithimage}

\item
  \begin{questionwithimage}{chuyi_xia_2-p024-cleaned-8.jpg}
  如图，取编号为1-6的6个由小三角形组成的图案中的5块恰好无空隙的填成左侧的大图案，图中显示的所有小三角形都是全等的正三角形，且每一个图案都可以任意旋转、翻转。6个图案中有一个是用不上的，其中用不上的那个图案是（\quad）
\circled{1}
\circled{2}
\circled{3}
\circled{4}
\circled{5}
  \optionlist
  \item (3)
  \item (2)
  \item (4)
  \item (5)
  \end{enumerate}
  \end{questionwithimage}

\item
  \begin{questionwithimage}{chuyi_xia_2-p024-cleaned-15.jpg}
  如图，矩形 ABCD 中，E 是 AD 的中点，将 $\triangle ABE$ 沿直线$BE$折叠后得到 $\triangle GBE$，延长 $BG$交 CD 于点$F$. 若 $AB=6, BC=4\sqrt{6}$，则 FD 的长为（\quad）
  \optionlist
  \item 2
  \item 4
  \item $\sqrt{6}$
  \item $2\sqrt{3}$
  \end{enumerate}
  \end{questionwithimage}

\item
  \begin{questionwithimage}{chuyi_xia_2-p024-cleaned-16.jpg}
  如图， $\triangle ABC$ 的顶点坐标分别为 $A(4, 6)$、 $B(5, 2)$、 $C(2, 1)$，如果将 $\triangle ABC$ 绕点 $C$ 按逆时针方向旋转 $90^\circ$，得到 $\triangle A'B'C$，将 $\triangle A'B'C$ 向下平移 5 个单位，得 $\triangle A''B''C''$，那么点 $A$ 的对应点 $A''$ 的坐标是（\quad）
  \optionlist
  \item $(-3,-2)$
  \item $(3,-8)$
  \item $(-2,-1)$
  \item $(1,-1)$
  \end{enumerate}
  \end{questionwithimage}

\end{enumerate}

\section*{p025}
\subsection*{四、未分类题目（共 2 小题）}

\solutionlist[7]
\item
  \begin{questionwithimage}{chuyi_xia_2-p025-cleaned-1.jpg}
  如图，Rt$\triangle ABC$中，$\angle$ACB=90$^\circ$，AC=3，BC=4，将边$AC$沿CE翻折，使点$A$落在AB上的点$D$处；再将边$BC$沿CF翻折，使点$B$落在CD的延长线上的点$B$'处，两条折痕与斜边$AB$分别交于点$E$、F，则线段$B$'F的长为（\quad）
  \optionlist
  \item $\frac{3}{5}$
  \item $\frac{4}{5}$
  \item $\frac{2}{3}$
  \item $\frac{\sqrt{3}}{2}$
  \end{enumerate}
  \end{questionwithimage}

\item
  坐标平面内有点 $A(4, 8)$ , $B(-4,-8)$ , 以坐标轴为对称轴, 点 $A$ 可以由点 $B$ 经过 $m$ 次轴对称变换得到, 则 $m$ 的最小值为 （\quad）
  \optionlist
  \item 1
  \item 2
  \item 3
  \item 4
  \end{enumerate}

\end{enumerate}

\subsection*{二、填空题（共 6 小题）}

\blanklist[9]
\item
  如图，在 Rt$\triangle ABC$中，$\angle$ACB=90$^\circ$，AC=8，BC=6，点$E$是 $AB$边上一动点，过点$E$作 $DE$$\perp$AB 交 $AC$边于点$D$，将 $\angle$A 沿直线$DE$翻折，点$A$落在线段$AB$上的 F 处，连接 $FC$，当 $\triangle BCF$为等腰三角形时，AE 的长为 \_\_\_\_.\underline{\hspace{2cm}}

  \centerquestionimage{0.35\textwidth}{chuyi_xia_2-p025-cleaned-2.jpg}

\item
  等边三角形 $ABC$ 在平面直角坐标系中， $A(1,0)$、 $B(4,0)$， $C$ 点在第一象限，点 $P$ 是 $\triangle ABC$ 三条高的交点，则点 $P$ 关于 $\triangle ABC$ 各边所在直线的对称点的坐标是 \_\_\_\_.\underline{\hspace{2cm}}

\item
  把一个三角形绕其中一个顶点逆时针旋转并放大或缩小（这个顶点不变），我们把这样的三角形运动称为三角形的 $T$ -变换，这个顶点称为 $T$ -变换中心，旋转角称为 $T$ -变换角，三角形与原三角形的对应边之比称为 $T$ -变换比；已知 $\triangle ABC$ 在直角坐标平面内，点 $A(0,-1)$， $B(-\sqrt{3}, 2)$， $C(0, 2)$，将 $\triangle ABC$ 进行 $T$ -变换， $T$ -变换中心为点 $A$， $T$ -变换角为 $60^{\circ}$， $T$ -变换比为 $\frac{2}{3}$，那么经过 $T$ -变换后点 $C$ 所对应的点的坐标为 \_\_\_\_.\underline{\hspace{2cm}}

\item
  在平面直角坐标系中，点 $A$ 的坐标为（2，0），点 $B$ 的坐标为（0，1），将线段 $AB$ 平移，使其一个端点到 $C(3,2)$，则平移后另一端点的坐标为 \_\_\_\_.\underline{\hspace{2cm}}

\item
  如图，四边形 $ABCD$ 是矩形纸片， $AB=2$。对折矩形纸片 $ABCD$，使 $AD$ 与 $BC$ 重合，折痕为 $EF$；展平后再过点 $B$ 折叠矩形纸片，使点 $A$ 落在 $EF$ 上的点 $N$，折痕 $BM$ 与 $EF$ 相交于点 $Q$；再次展平，连接 $BN$， $MN$，延长 $MN$ 交 $BC$ 于点 $G$。有如下结论：
\circled{1} $\angle ABN=60^{\circ}$；\circled{2} $AM=1$；\circled{3} $QN=\frac{\sqrt{3}}{3}$；\circled{4} $\triangle BMG$ 是等边三角形；\circled{5} $PN+PH=\sqrt{3}$。其中正确结论的序号是 \_\_\_\_。\underline{\hspace{2cm}}

  \centerquestionimage{0.35\textwidth}{chuyi_xia_2-p025-cleaned-3.jpg}

\item
  如图，在平面直角坐标系中，O 为原点，点$A$ 的坐标为 $(-4, 0)$，点$B$ 的坐标为 $(0, 4)$，点$C$、D 分别为 OA、OB 的中点，若正方形 OCED 绕点$O$顺时针旋转，得正方形 $OC'E'D'$。记旋转角为 $a(0^{\circ} \leqslant a < 360^{\circ})$，连接 $AC'$、 $BD'$，设直线 $AC'$ 与直线 $BD'$ 相交于点$F$，则点$F$的纵坐标的最大值为 \_\_\_\_。\underline{\hspace{2cm}}

  \centerquestionimage{0.35\textwidth}{chuyi_xia_2-p025-cleaned-4.jpg}

\end{enumerate}

\section*{p026}
\subsection*{三、解答题（共 2 小题）}

\solutionlist[15]
\item
  如图\circled{1}，正方形ABCD与正方形AEFG的边 $AB$、 $AE$ （$AB < AE$）在一条直线上，正方形AEFG以点 $A$ 为旋转中心逆时针旋转，设旋转角为 $\alpha$。在旋转过程中，两个正方形只有点 $A$ 重合，其它顶点均不重合，连接 $BE$、 $DG$
（1）当正方形AEFG旋转至如图\circled{2}所示的位置时，求证：BE=DG;
(2) 如图\circled{3}, 如果 $\alpha=45^{\circ}, AB=2, AE=4\sqrt{2}$ , 求点 $G$ 到 $BE$ 的距离.

  \begin{questionimagegroup}
  \questionimageitem{0.28\textwidth}{chuyi_xia_2-p026-cleaned-1.jpg}{图\circled{1}}
  \questionimagegap
  \questionimageitem{0.28\textwidth}{chuyi_xia_2-p026-cleaned-2.jpg}{图\circled{2}}
  \questionimagegap
  \questionimageitem{0.28\textwidth}{chuyi_xia_2-p026-cleaned-3.jpg}{图\circled{3}}
  \end{questionimagegroup}

\item
  在等边三角形 $ABC$ 中, $AD \perp BC$ 于点 $D$ .
(1) 如图\circled{1}, 请你直接写出线段 ${AD}$ 与 ${BC}$ 之间的数量关系： ${AD}=$ \_\_\_\_\_ ${BC}$；
（2）如图\circled{2}，若 $P$ 是线段 $BC$ 上一个动点（点 $P$ 不与点 $B$、 $C$ 重合），连接 $AP$，将线段 $AP$ 绕点 $A$ 逆时针旋转 $60^{\circ}$，得到线段 $AE$，连接 $CE$，猜想线段 $AD$、 $CE$、 $PC$ 之间的数量关系，并证明你的结论；
（3）如图\circled{3}，若点 $P$ 是线段 $BC$ 延长线上一个动点，（2）中的其他条件不变，按照（2）中的作法，请在图\circled{3}中补全图形，并直接写出线段 $AD$、 $CE$、 $PC$ 之间的数量关系.

  \begin{questionimagegroup}
  \questionimageitem{0.28\textwidth}{chuyi_xia_2-p026-cleaned-4.jpg}{图\circled{1}}
  \questionimagegap
  \questionimageitem{0.28\textwidth}{chuyi_xia_2-p026-cleaned-5.jpg}{图\circled{2}}
  \questionimagegap
  \questionimageitem{0.28\textwidth}{chuyi_xia_2-p026-cleaned-6.jpg}{图\circled{3}}
  \questionimagegap
  \questionimageitem{0.28\textwidth}{chuyi_xia_2-p026-cleaned-7.jpg}{}
  \end{questionimagegroup}

\end{enumerate}

\section*{p027}
\subsection*{四、未分类题目（共 2 小题）}

\solutionlist[17]
\item
  把一副三角板如图甲放置，其中 $\angle ACB=\angle DEC=90^{\circ}$， $\angle A=45^{\circ}$， $\angle D=30^{\circ}$， $BC=\sqrt{3}+1$。斜边 $AB$、 $DC$ 相交于点 $O$。
(1) 求 CO 的长;
（2）若把三角板 $DCE$ 绕点 $C$ 顺时针旋转 $15^{\circ}$ 得到 $\triangle D_{1}CE_{1}$ （如图乙），这时 $AB$ 与 $CD_{1}$ 相交于点 $O_{1}$，此时，求： $CO_{1}$ 的长；
（3）若把三角板 $D_{1}CE_{1}$ 绕着点 $C$ 顺时针再旋转 $15^{\circ}$ 得 $\triangle D_{2}CE_{2}$ （如图丙），这时 $AB$ 与 $CD_{2}$ 相交于点 $O_2$，此时，求： $CO_{2}$ 的长.

  \begin{questionimagegroup}
  \questionimageitem{0.28\textwidth}{chuyi_xia_2-p027-cleaned-1.jpg}{(甲)}
  \questionimagegap
  \questionimageitem{0.28\textwidth}{chuyi_xia_2-p027-cleaned-2.jpg}{(乙)}
  \questionimagegap
  \questionimageitem{0.28\textwidth}{chuyi_xia_2-p027-cleaned-3.jpg}{}
  \end{questionimagegroup}

\item
  如图，在 Rt$\triangle ABC$中，$\angle$C=90$^\circ$，点$P$为 $AC$边上的一点，将线段$AP$绕点$A$顺时针方向旋转（点$P$对应点$P$'），当 AP 旋转至 AP'$\perp$AB 时，点$B$、P、P' 恰好在同一直线上，此时作 P'E$\perp$AC 于点$E$.
(1) 求证: $\angle CBP=\angle ABP;$
(2) 若 AB-BC=4, AC=8, 求 AE 的长;
（3）当 $\angle ABC=60^{\circ}$， $BC=2$ 时，点 $N$ 为 $BC$ 的中点，点 $M$ 为边 $BP$ 上一个动点，连接 $MC$， $MN$ 求 $MC+MN$ 的最小值.

  \begin{questionimagegroup}
  \questionimageitem{0.32\textwidth}{chuyi_xia_2-p027-cleaned-4.jpg}{}
  \questionimagegap
  \questionimageitem{0.32\textwidth}{chuyi_xia_2-p027-cleaned-5.jpg}{}
  \end{questionimagegroup}

\end{enumerate}

\section*{p028}
\subsection*{一、选择题（共 7 小题）}

\questionlist
\item
  \begin{questionwithimage}{chuyi_xia_2-p028-cleaned-1.jpg}
  如图，在 $\square ABCD$ 中，用直尺和圆规作 $\angle BAD$ 的平分线 $AG$ 交 $BC$ 于点 $E$ . 若 $BF=6, AB=5$，则 $AE$ 的长为（\quad）
  \optionlist
  \item 4
  \item 6
  \item 8
  \item 10
  \end{enumerate}
  \end{questionwithimage}

\item
  如图，在 $\square ABCD$ 中， $\angle A=70^{\circ}$，将 $\square ABCD$ 折叠，使点 $D$、 $C$ 分别落在点 $F$、 $E$ 处（点 $F$、 $E$ 都在 $AB$ 所在的直线上），折痕为 $MN$，则 $\angle AMF$ 等于（\quad）
  \optionlist
  \item $70^{\circ}$
  \item $40^{\circ}$
  \item $30^{\circ}$
  \item $20^{\circ}$
  \end{enumerate}

\item
  \begin{questionwithimage}{chuyi_xia_2-p028-cleaned-2.jpg}
  平行四边形的一条边长是 $12cm$ , 那么它的两条对角线的长可能是 （\quad）
  \optionlist
  \item $8cm$ 和 $16cm$
  \item $10cm$ 和 $16cm$
  \item $8cm$ 和 $14cm$
  \item $8cm$ 和 $12cm$
  \end{enumerate}
  \end{questionwithimage}

\item
  \begin{questionwithimage}{chuyi_xia_2-p028-cleaned-3.jpg}
  如图，以 $\square ABCD$ 的边 $CD$ 为斜边向内作等腰直角 $\triangle CDE$，使 $AD=DE=CE$， $\angle DEC=90^\circ$，且点 $E$ 在平行四边形内部，连接 $AE$、 $BE$，则 $\angle AEB$ 的度数是（\quad）
  \optionlist
  \item $120^\circ$
  \item $135^\circ$
  \item $150^\circ$
  \item $45^\circ$
  \end{enumerate}
  \end{questionwithimage}

\item
  \begin{questionwithimage}{chuyi_xia_2-p028-cleaned-4.jpg}
  如图，平面直角坐标系中，平行四边形 ABCD 的中心 E 的坐标为 $(2, 0)$，若点$A$ 的坐标为 $(-2, 1)$，则点$C$ 的坐标为（\quad）
  \optionlist
  \item $(4,-1)$
  \item $(6,-1)$
  \item $(8,-1)$
  \item $(6,-2)$
  \end{enumerate}
  \end{questionwithimage}

\item
  在面积为 60 的 $\square$ABCD 中，过点$A$作 $AE \perp $ 直线$BC$于点$E$，作 $AF \perp $ 直线$CD$于点$F$，若 AB = 10，BC = 12，则 $CE+CF$ 的值为（\quad）
  \optionlist
  \item $22+11\sqrt{3}$
  \item $22-11\sqrt{3}$
  \item $22+11\sqrt{3}$ 或 $22-11\sqrt{3}$
  \item $22+11\sqrt{3}$ 或 $2+\sqrt{3}$
  \end{enumerate}

\item
  \begin{questionwithimage}{chuyi_xia_2-p028-cleaned-5.jpg}
  如图，在 $\square ABCD$ 中，对角线 $AC$ 与 $BD$ 相交于点 $E$， $\angle AEB=45^{\circ}$， $BD=2$，将 $\triangle ABC$ 沿 $AC$ 所在直线翻折 $180^{\circ}$ 到其原来所在的同一平面内，若点 $B$ 的落点记为 $B'$，则 $DB'$ 的长为（\quad）
  \optionlist
  \item $\sqrt{2}$
  \item $\sqrt{3}$
  \item $2\sqrt{2}$
  \item $\sqrt{5}$
  \end{enumerate}

  % 图注：图\circled{1}
  \end{questionwithimage}

\end{enumerate}

\section*{p029}
\subsection*{四、未分类题目（共 1 小题）}

\solutionlist[8]
\item
  如图，在 $\square ABCD$ 中， $AB=3$， $AD=4$， $\angle ABC=60^{\circ}$，过 $BC$ 的中点 $E$ 作 $EF \perp AB$，垂足为点 $F$，与 $DC$ 的延长线相交于点 $H$，则 $\triangle DEF$ 的面积是（\quad）
  \optionlist
  \item $2\sqrt{3}$
  \item $4\sqrt{3}$
  \item $3+\sqrt{3}$
  \item $6+2\sqrt{3}$
  \end{enumerate}

\end{enumerate}

\subsection*{二、填空题（共 6 小题）}

\blanklist[9]
\item
  在 $\square ABCD$ 中， $E, F$ 是对角线 $AC$ 上两点，要使 $BE=DF$，须添上一个适当的条件是：\_\_\_\_。（填一个即可）\underline{\hspace{2cm}}

\item
  在直角坐标系中，已知 $A(1,0)$、 $B(-1,-2)$、 $C(2,-2)$ 三点坐标，若以 $A$、 $B$、 $C$、 $D$ 为顶点的四边形是平行四边形，那么点 $D$ 的坐标可以是 \_\_\_\_.
\circled{1} $(-2, 0)$ \circled{2} $(0,-4)$ \circled{3} $(4, 0)$ \circled{4} $(1,-4)$\underline{\hspace{2cm}}

\item
  如图， $\square ABCD$ 中，点 $E$ 在边 $AD$ 上，以 $BE$ 为折痕，将 $\triangle ABE$ 折叠，使点 $A$ 正好落在 $CD$ 上的点 $F$，若 $\triangle FDE$ 的周长为 8， $\triangle FCB$ 的周长为 14，则 $FC$ 的长为 \_\_\_\_.\underline{\hspace{2cm}}

\item
  如图， $\triangle ACE$ 是以 $\square ABCD$ 的对角线 $AC$ 为边的等边三角形，点 $C$ 与点 $E$ 关于 $x$ 轴对称。若 $E$ 点的坐标是（7， $-3\sqrt{3}$），则 $D$ 点的坐标是 \_\_\_\_.\underline{\hspace{2cm}}

\item
  如图，在平行四边形 ABCD 中，分别以 AB、AD 为边向外作等边 $\triangle ABE$、 $\triangle ADF$，延长 $CB$交 AE 于点$G$，点$G$在点$A$，E 之间，连接 $CE$、CF、EF，有下列四个结论：
\circled{1} $\triangle CDF\cong \triangle EBC$；\circled{2} $\angle CDF=\angle EAF$；\circled{3} $\triangle ECF$ 是等边三角形；\circled{4} $CG \perp AE$。其中正确的结论的个数是\_\_\_\_个。\underline{\hspace{2cm}}

\item
  如图，对面积为1的$\square$ABCD逐次进行以下操作：第一次操作，分别延长 $AB$，BC，CD，DA至点 $A_{1}$， $B_{1}$， $C_{1}$， $D_{1}$，使得 $A_{1}B=2AB$， $B_{1}C=2BC$， $C_{1}D=2CD$， $D_{1}A=2AD$，顺次连接 $A_{1}$， $B_{1}$， $C_{1}$， $D_{1}$，得到$\square$ $A_{1}B_{1}C_{1}D_{1}$，记其面积为 $S_{1}$；第二次操作，分别延长 $A_{1}B_{1}$， $B_{1}C_{1}$， $C_{1}D_{1}$、 $D_{1}A_{1}$ 至点 $A_{2}$， $B_{2}$， $C_{2}$， $D_{2}$，使得 $A_{2}B_{1}=2A_{1}B_{1}$， $B_{2}C_{1}=2B_{1}C_{1}$， $C_{2}D_{1}=2C_{1}D_{1}$， $D_{2}A_{1}=2A_{1}D_{1}$，顺次连接 $A_{2}$， $B_{2}$， $C_{2}$， $D_{2}$ 记其面积为 $S_{2}$；…；按此规律继续下去，可得到$\square$ $A_{5}B_{5}C_{5}D_{5}$，则其面积 $S_{5}=$ \_\_\_\_。\underline{\hspace{2cm}}

  \begin{questionimagegroup}
  \questionimageitem{0.28\textwidth}{chuyi_xia_2-p029-cleaned-1.jpg}{}
  \questionimagegap
  \questionimageitem{0.28\textwidth}{chuyi_xia_2-p029-cleaned-2.jpg}{}
  \questionimagegap
  \questionimageitem{0.28\textwidth}{chuyi_xia_2-p029-cleaned-3.jpg}{}
  \questionimagegap
  \questionimageitem{0.28\textwidth}{chuyi_xia_2-p029-cleaned-4.jpg}{}
  \questionimagegap
  \questionimageitem{0.28\textwidth}{chuyi_xia_2-p029-cleaned-5.jpg}{}
  \questionimagegap
  \questionimageitem{0.28\textwidth}{chuyi_xia_2-p029-cleaned-6.jpg}{}
  \questionimagegap
  \questionimageitem{0.28\textwidth}{chuyi_xia_2-p029-cleaned-7.jpg}{}
  \end{questionimagegroup}

\end{enumerate}

\section*{p030}
\subsection*{三、解答题（共 3 小题）}

\solutionlist[15]
\item
  如图，已知 $\square ABCD$ 中， $DE \perp BC$ 于点 $E$， $DH \perp AB$ 于点 $H$， $AF$ 平分 $\angle BAD$，分别交 $DC$、 $DE$、 $DH$ 于点 $F$、 $G$、 $M$，且 $DE=AD$ .
(1) 求证: $\triangle ADG\cong \triangle FDM$ .
(2) 猜想 AB 与 $DG+CE$ 之间有何数量关系，并证明你的猜想.

  \centerquestionimage{0.35\textwidth}{chuyi_xia_2-p030-cleaned-1.jpg}

\item
  在 $\square ABCD$ 中，对角线 $BD \perp BC$， $G$ 为 $BD$ 延长线上一点且 $\triangle ABG$ 为等边三角形， $\angle BAD$、 $\angle CBD$ 的平分线相交于点 $E$，连接 $AE$ 交 $BD$ 于 $F$，连接 $GE$。若平行四边形 $ABCD$ 的面积为 $9\sqrt{3}$，求 $AG$ 的长。

  \centerquestionimage{0.35\textwidth}{chuyi_xia_2-p030-cleaned-2.jpg}

\item
  已知如图，四边形 $ABCD$ 为平行四边形， $AD=a$， $AC$ 为对角线， $BM \parallel AC$，过点 $D$ 作 $DE \parallel CM$，交 $AC$ 的延长线于 $F$，交 $BM$ 的延长线于 $E$ .
(1) 求证: $\triangle ADF\cong \triangle BCM;$
(2) 若 ${AC}={2CF}$ , $\angle {ADC}={60}^{\circ}$ , ${AC} \perp {DC}$ ,求四边形 ${ABED}$ 的面积(用含 $a$ 的代数式表示).

  \centerquestionimage{0.35\textwidth}{chuyi_xia_2-p030-cleaned-3.jpg}

\end{enumerate}

\section*{p031}
\subsection*{一、选择题（共 7 小题）}

\questionlist
\item
  \begin{questionwithimage}{chuyi_xia_2-p031-cleaned-1.jpg}
  在平面直角坐标系中；以 $A(1,1), B(3,0), C(-1,0)$ 为顶点构造平行四边形，下列各点不能作为平行四边形顶点的是（\quad）
  \optionlist
  \item (5, 1)
  \item (0, -2)
  \item (-3, 1)
  \item (1, -1)
  \end{enumerate}
  \end{questionwithimage}

\item
  \begin{questionwithimage}{chuyi_xia_2-p031-cleaned-2.jpg}
  如图, 在四边形 $ABCD$ 中, $AC$ 与 $BD$ 相交于点 $O$ , 如果只给出条件 “ $AB//CD$ ”, 还不能判定四边形 $ABCD$ 为平行四边形, 若想使四边形 $ABCD$ 为平行四边形, 要添加一个条件, 这个条件可以是 （\quad） \circled{1}如果再添加条件 “ $BC=AD$ ” \circled{2}如果再添加条件 “ $\angle BAD=\angle BCD$ ”, \circled{3}如果再添加条件 “ $OA=OC$ ” \circled{4}如果再添加条件 “ $\angle ABD=\angle CAB$ .
  \optionlist
  \item \circled{1}或\circled{2}
  \item \circled{1}或\circled{3}或\circled{4}
  \item \circled{2}或\circled{3}
  \item \circled{2}或\circled{3}或\circled{4}
  \end{enumerate}
  \end{questionwithimage}

\item
  下列说法中，不正确是（\quad）
  \optionlist
  \item 对角线互相平分的四边形是平行四边形
  \item 两组对角分别相等的四边形是平行四边形
  \item 一组对边平行且相等的四边形是平行四边形
  \item 一组对边平行另一组对边相等的四边形是平行四边形
  \end{enumerate}

\item
  下列选项中的四边形只有一个为平行四边形，根据图中所给的边长长度及角度，哪一个图形为平行四边形？（）

  \begin{questionimagegroup}
  \questionimageitem{0.28\textwidth}{chuyi_xia_2-p031-cleaned-3.jpg}{}
  \questionimagegap
  \questionimageitem{0.28\textwidth}{chuyi_xia_2-p031-cleaned-4.jpg}{}
  \questionimagegap
  \questionimageitem{0.28\textwidth}{chuyi_xia_2-p031-cleaned-5.jpg}{}
  \questionimagegap
  \questionimageitem{0.28\textwidth}{chuyi_xia_2-p031-cleaned-6.jpg}{}
  \end{questionimagegroup}

\item
  点$A$, B, C, D 在同一平面内，从四个条件中（1）AB = CD，（2）AB // CD，（3）BC = AD，（4）BC // AD 中任选两个，使四边形 ABCD 是平行四边形，这样的选法有（\quad）
  \optionlist
  \item 3 种
  \item 4 种
  \item 5 种
  \item 6 种
  \end{enumerate}

\item
  在下列给定的条件中，能画出平行四边形的是（\quad）
  \optionlist
  \item 以 60cm 为一条对角线，20cm，34cm 为两邻边
  \item 以 6cm，10cm 为对角线，8cm 为一边
  \item 以 6cm 为一条对角线，3cm 和 10cm 为两邻边
  \item 以 20cm，36cm 为对角线，22cm 为一边
  \end{enumerate}

\item
  在四边形 $ABCD$ 中， $AC$ 与 $BD$ 相交于点 $O$，如果只给出条件“ $AB \parallel CD$ ”，那么还不能判定四边形 $ABCD$ 为平行四边形，给出以下 5 个说法：
\circled{1}如果再加上条件“ $\angle DBA=\angle CAB$ ”，那么四边形 $ABCD$ 为平行四边形。
\circled{2}如果再加上条件“ $AB=CD$ ”，那么四边形 $ABCD$ 为平行四边形；
\circled{3}如果再加上条件“ $\angle DAB=\angle DCB$ ”，那么四边形 $ABCD$ 为平行四边形；
\circled{4}如果再加上条件“ $BC=AD$ ”，那么四边形 $ABCD$ 是平行四边形；
\circled{5}如果再加上条件“ $AO=CO$ ”，那么四边形 $ABCD$ 为平行四边形；
其中正确的说法有（\quad）
  \optionlist
  \item 5个
  \item 4个
  \item 3个
  \item 2个
  \end{enumerate}

\end{enumerate}

\section*{p032}
\subsection*{四、未分类题目（共 1 小题）}

\solutionlist[8]
\item
  \begin{questionwithimage}{chuyi_xia_2-p032-cleaned-1.jpg}
  如图，甲、乙两人想在正五边形 $ABCDE$ 内部找一点 $P$，使得四边形 $ABPE$ 为平行四边形，其作法如下：
（甲）连接 $BD$、 $CE$，两线段相交于 $P$ 点，则 $P$ 即为所求
（乙）先取 $CD$ 的中点 $M$，再以 $A$ 为圆心， $AB$ 长为半径画弧，交 $AM$ 于 $P$ 点，则 $P$ 即为所求.
对于甲、乙两人的作法，下列判断正确的是（\quad）
  \optionlist
  \item 两人都正确
  \item 两人都错误
  \item 甲正确，乙错误
  \item 甲错误，乙正确
  \end{enumerate}
  \end{questionwithimage}

\end{enumerate}

\subsection*{二、填空题（共 6 小题）}

\blanklist[9]
\item
  如图，在矩形 ABCD 中，点$E$，F 分别在边$AB$，CD 上，请添加一个条件 \_\_\_\_，使四边形 AECF 是平行四边形（只填一个即可）.\underline{\hspace{2cm}}

  \centerquestionimage{0.35\textwidth}{chuyi_xia_2-p032-cleaned-2.jpg}

\item
  在平面直角坐标系中，已知 $A(-2, 1), B(-2,-1), O(0, 0)$。若以 $A, B, C, O$ 为顶点的四边形为平行四边形，那么点 $C$ 的坐标是 \_\_\_\_.\underline{\hspace{2cm}}

\item
  如图所示，在格点图中，以格点 $A$、$B$、$C$、$D$、$E$、$F$ 为顶点，你能画出 \_\_\_\_ 个平行四边形，并在图中画出来。\underline{\hspace{2cm}}

  \centerquestionimage{0.35\textwidth}{chuyi_xia_2-p032-cleaned-3.jpg}

\item
  如图，在 $Rt\triangle ABC$ 中， $\angle BAC=90^{\circ}$， $AB=3$， $AC=4$，将 $\triangle ABC$ 沿直线 $BC$ 向右平移 2.5 个单位得到 $\triangle DEF$， $AC$ 与 $DE$ 相交于 $G$ 点，连接 $AD$， $AE$，则下列结论中成立的是 \_\_\_\_。
\circled{1}四边形 $ABED$ 是平行四边形；\circled{2} $\triangle AGD\cong \triangle CGE$；
\circled{3} $\triangle ADE$ 为等腰三角形；\circled{4} $AC$ 平分 $\angle EAD$。\underline{\hspace{2cm}}

  \centerquestionimage{0.35\textwidth}{chuyi_xia_2-p032-cleaned-4.jpg}

\item
  在平面直角坐标系 $xOy$ 中，有 $A(3,2), B(-1,-4)$， $P$ 是 $x$ 轴上的一点， $Q$ 是 $y$ 轴上的一点，若以点 $A, B, P, Q$ 四个点为顶点的四边形是平行四边形，则 $Q$ 点的坐标是 \_\_\_\_ \_\_\_\_ \_\_\_\_ \_\_\_\_ \_\_\_\_ \_\_\_\_ \_\_\_\_ \_\_\_\_ \_\_\_\_ \_\_\_\_ \_\_\_\_ \_\_\_\_ \_\_\_\_ \_\_\_\_ \_\_\_\_ \_\_\_\_ \_\_\_\_ \_\_\_\_ \_\_\_\_ \_\_\_\_ \_\_\_\_ \_\_\_\_ \_\_\_\_ \_\_\_\_ \_\_\_\_ \_\_\_\_ \_\_\_\_ \_\_\_\_ \_\_\_\_ \_\_\_\_ \_\_\_\_ \_\_\_\_ \_\_\_\_ \_\_\_\_ \_\_\_\_ \_\_\_\_ \_\_\_\_ \_\_\_\_ \_\_\_\_ \_\_\_\_ \_\_\_\_ \_\_\_\_ \_\_\_\_ \_\_\_\_ \_\_\_\_ \_\_\_\_ \_\_\_\_ \_\_\_\_ \_\_\_\_ \_\_\_\_ \_\_\_\_ \_\_\_\_ \_\_\_\_ \_\_\_\_ \_\_\_\_ \_\_\_\_ \_\_\_\_ \_\_\_\_ \_\_\_\_ \_\_\_\_ \_\_\_\_ \_\_\_\_ \_\_\_\_ \_\_\_\_ \_\_\_\_ \_\_\_\_ \_\_\_\_ \_\_\_\_ \_\_\_\_ \_\_\_\_ \_\_\_\_ \_\_\_\_ \_\_\_\_ \_\_\_\_ \_\_\_\_ \_\_\_\_ \_\_\_\_ \_\_\_\_ \_\_\_\_ \_\_\_\_ \_\_\_\_ \_\_\_\_ \_\_\_\_ \_\_\_\_ \_\_\_\_ \_\_\_\_ \_\_\_\_ \_\_\_\_ \_\_\_\_ \_\_\_\_ \_\_\_\_ \_\_\_\_ \_\_\_\_ \_\_\_\_ \_\_\_\_ \_\_\_\_ \_\_\_\_ \_\_\_\_ \_\_\_\_ \_\_\_\_ \_\_\_\_.\underline{\hspace{2cm}}

\item
  如图，在等边三角形 $ABC$ 中， $BC=6cm$，射线 $AG // BC$，点 $E$ 从点 $A$ 出发沿射线 $AG$ 以 $1\mathrm{cm/s}$ 的速度运动，点 $F$ 从点 $B$ 出发沿射线 $BC$ 以 $2\mathrm{cm/s}$ 的速度运动。如果点 $E$、 $F$ 同时出发，设运动时间为 $t$ （s）当 $t=\_\_\_\_$ s 时，以 $A$、 $C$、 $E$、 $F$ 为顶点四边形是平行四边形。\underline{\hspace{2cm}}

  \centerquestionimage{0.35\textwidth}{chuyi_xia_2-p032-cleaned-5.jpg}

\end{enumerate}

\subsection*{三、解答题（共 1 小题）}

\solutionlist[15]
\item
  我们给出如下定义：若一个四边形的两条对角线相等，则称这个四边形为等对角线四边形。请解答下列问题：
（1）写出你所学过的特殊四边形中是等对角线四边形的两种图形的名称；
（2）探究：当等对角线四边形中两条对角线所夹锐角为 $60^{\circ}$ 时，这对 $60^{\circ}$ 角所对的两边之和与其中一条对角线的大小关系，并证明你的结论.

  \centerquestionimage{0.35\textwidth}{chuyi_xia_2-p032-cleaned-6.jpg}

\end{enumerate}

\section*{p033}
\subsection*{四、未分类题目（共 2 小题）}

\solutionlist[16]
\item
  以 $\triangle ABC$ 的边 $AB$、 $AC$ 为直角边向外作等腰直角 $\triangle ABD$ 和等腰直角 $\triangle ACE$， $M$ 是 $BC$ 中点，连接 $AM$ 和 $DE$ .
（1）如图\circled{1}， $\triangle ABC$ 中 $\angle BAC=90^{\circ}$ 时， $AM$ 与 $ED$ 大小的关系是 \_\_\_\_。 $AM$ 与 $ED$ 的位置关系是 \_\_\_\_；
(2) 如图\circled{2}, $\triangle ABC$ 为一般三角形时线段 $AM$ 与 $ED$ 的关系是\_\_\_\_. 试证明你的结论;
（3）如图\circled{3}，若以 $\triangle ABC$ 的边 $AB$、 $AC$ 为直角边，向内作等腰直角 $\triangle ABE$ 和 $\triangle ACD$，其它条件不变，试探究线段 $AM$ 与 $DE$ 之间的关系，不要求证明你的结论.

  \begin{questionimagegroup}
  \questionimageitem{0.28\textwidth}{chuyi_xia_2-p033-cleaned-1.jpg}{图\circled{1}}
  \questionimagegap
  \questionimageitem{0.28\textwidth}{chuyi_xia_2-p033-cleaned-2.jpg}{}
  \questionimagegap
  \questionimageitem{0.28\textwidth}{chuyi_xia_2-p033-cleaned-3.jpg}{图\circled{3}}
  \end{questionimagegroup}

\item
  $\triangle ABC$ 是等边三角形，点 $D$ 是射线上 $BC$ 上的一个动点（点 $D$ 不与点 $B$， $C$ 重合， $\triangle ADE$ 是以 $AD$ 为边的等边三角形，过点 $E$ 作 $BC$ 的平行线，分别交射线 $AB$， $AC$ 于点 $F$， $G$，连接 $BE$
(1) 如图\circled{1} 所示, 当点$D$在线段$BC$上时:
\circled{1}求证： $\triangle AEB\cong \triangle ADC$；\circled{2}探究四边形 $BCGE$ 是哪种特殊的四边形，并说明理由.
探究四边形 $BCGE$ 是哪种特殊的四边形，并说明理由。如图\circled{2}所示，当点 $D$ 在 $BC$ 的延长线上时，直接写出（1）中的两个结论是否成立。

  \begin{questionimagegroup}
  \questionimageitem{0.32\textwidth}{chuyi_xia_2-p033-cleaned-4.jpg}{图\circled{1}}
  \questionimagegap
  \questionimageitem{0.32\textwidth}{chuyi_xia_2-p033-cleaned-5.jpg}{图\circled{2}}
  \end{questionimagegroup}

\end{enumerate}

\section*{p034}
\subsection*{四、未分类题目（共 1 小题）}

\solutionlist[18]
\item
  两个全等的三角形如下图所示放置，点 $B$、 $A$、 $D$ 在同一直线上．操作：在图中，在 $CB$ 边上截取 $GM=AB$，连接 $DM$，交 $AC$ 于 $N$ ．请探究$\angle$AND的大小，并证明你的结论.

  \centerquestionimage{0.35\textwidth}{chuyi_xia_2-p034-cleaned-1.jpg}

\end{enumerate}

\section*{p035}
\subsection*{一、选择题（共 6 小题）}

\questionlist
\item
  如图，已知在 $\square ABCD$ 中， $E, F$ 是对角线 $BD$ 上的两点，则以下条件不能判断四边形 $AECF$ 为平行四边形的是（\quad）
  \optionlist
  \item $BE=DF$
  \item $AF \perp BD, CE \perp BD$
  \item $\angle BAE=\angle DCF$
  \item $AF=CE$
  \end{enumerate}

\item
  如图所示，在 $\square ABCD$ 中， $E, F$ 分别是 $AB, CD$ 的中点，连接 $DE, EF, BF$，则图中平行四边形共有（\quad）
  \optionlist
  \item 2个
  \item 4个
  \item 6个
  \item 8个
  \end{enumerate}

\item
  如图，已知 $\square$ABCD 中，$\angle$DBC=45$^\circ$，DE$\perp$BC 于 E，BF$\perp$CD 于 F，D（\quad）
  \optionlist
  \item BF 交于 H，B
  \item AD 的延长线交于 G，下面结论正确的是（）
\circled{1} $DB=\sqrt{2}BE$ ; \circled{2} $\angle A=\angle BHE$ ; \circled{3} 域 CG，则四边形 BCGD 为平行四边形；\circled{4} $AD^{2}+DH^{2}=2DC^{2}$ .
  \item \circled{1}\circled{2}\circled{3}\circled{4}
  \item \circled{1}\circled{2}\circled{3}
  \item \circled{1}\circled{2}\circled{4}
  \item \circled{2}\circled{3}\circled{4}
  \end{enumerate}

\item
  如图，点 $D$ 是 $\triangle ABC$ 的边 $AB$ 的延长线上一点，点 $F$ 是边 $BC$ 上的一个动点（不与点 $B$ 重合）。以 $BD$、 $BF$ 为邻边作平行四边形 $BDEF$，又 $AP=BE$ 且 $AP \parallel BE$ （点 $P$、 $E$ 在直线 $AB$ 的同侧），如果 $BD=\frac{1}{4} AB$，那么 $\triangle PBC$ 的面积与 $\triangle ABC$ 面积之比为（\quad）
  \optionlist
  \item $\frac{1}{4}$
  \item $\frac{3}{5}$
  \item $\frac{1}{5}$
  \item $\frac{3}{4}$
  \end{enumerate}

\item
  如图，平行四边形 ABCD 中， $\angle ABD=30^{\circ}$，AB=4， $AE \perp BD$， $CF \perp BD$，且（\quad）
  \optionlist
  \item F 恰好是 BD 的三等分点，又 M、N 分别是 AB，CD 的中点，那么四边形 MENF 的面积是（）
  \item $3\sqrt{3}$
  \item $\sqrt{3}$
  \item $3\sqrt{2}$
  \item $\sqrt{2}$
  \end{enumerate}

\item
  \begin{questionwithimage}{chuyi_xia_2-p035-cleaned-1.jpg}
  如图， $\triangle ABC$ 是等边三角形， $P$ 是形内一点， $PD \parallel AB$， $PE \parallel BC$， $PF \parallel AC$，若 $\triangle ABC$ 的周长为 18，则 $PD+PE+PF=（\quad）$
  \optionlist
  \item 18
  \item $9\sqrt{3}$
  \item 6
  \item 条件不够，不能确定
  \end{enumerate}
  \end{questionwithimage}

\end{enumerate}

\section*{p036}
\subsection*{四、未分类题目（共 2 小题）}

\solutionlist[7]
\item
  \begin{questionwithimage}{chuyi_xia_2-p036-cleaned-1.jpg}
  如图,在平行四边形 $ABCD$ 中,按下列条件得到的四边形 $EFGH$ 不一定是平行四边形的是 （\quad）
  \optionlist
  \item 
  \item 
  \item 
  \end{enumerate}
  \end{questionwithimage}

\item
  \begin{questionwithimage}{chuyi_xia_2-p036-cleaned-5.jpg}
  如图，平行四边形 $ABCD$ 中， $AB=8cm$， $AD=12cm$，点 $P$ 在 $AD$ 边上以每秒 $1cm$ 的速度从点 $A$ 向点 $D$ 运动，点 $Q$ 在 $BC$ 边上，以每秒 $4cm$ 的速度从点 $C$ 出发，在 $CB$ 间往返运动，两个点同时出发，当点 $P$ 到达点 $D$ 时停止（同时点 $Q$ 也停止），在运动以后，以 $P$、 $D$、 $Q$、 $B$ 四点组成平行四边形的次数有（\quad）
  \optionlist
  \item 4次
  \item 3次
  \item 2次
  \item 1次
  \end{enumerate}
  \end{questionwithimage}

\end{enumerate}

\subsection*{二、填空题（共 6 小题）}

\blanklist[9]
\item
  如图， $\square ABCD$ 中， $\angle ABC=60^{\circ}$， $E$、 $F$ 分别在 $CD$ 和 $BC$ 的延长线上， $AE \parallel BD$， $EF \perp BC$， $EF=\div$ 则 $AB$ 的长是 \_\_\_\_.\underline{\hspace{2cm}}

  \centerquestionimage{0.35\textwidth}{chuyi_xia_2-p036-cleaned-6.jpg}

\item
  如图，在 $\triangle ABC$ 中， $AB=4$， $AC=3$， $BC=5$， $\triangle ABD$、 $\triangle ACE$、 $\triangle BCF$ 都是等边三角形，则四边形 $AEFD$ 的面积为\_\_\_\_.\underline{\hspace{2cm}}

  \centerquestionimage{0.35\textwidth}{chuyi_xia_2-p036-cleaned-7.jpg}

\item
  如图所示，六边形 $ABCDEF$ 中， $AB$ 平行且等于 $ED$， $AF$ 平行且等于 $CD$， $BC$ 平行且等于 $FE$，对角线 $FD \perp BD$。已知 $FD=24cm$， $BD=18cm$。则六边形 $ABCDEF$ 的面积是\_\_\_\_平方厘米。\underline{\hspace{2cm}}

  \centerquestionimage{0.35\textwidth}{chuyi_xia_2-p036-cleaned-8.jpg}

\item
  如图，在 $\triangle ABC$ 中， $\angle B$、 $\angle C$ 的平分线相交于 $O$ 点。作 $MN//BC$， $EF//AB$， $GH//AC$， $BC=a$， $AC=b$， $AB=c$，则 $\triangle GMO$ 周长+ $\triangle ENO$ 的周长- $\triangle FHO$ 的周长为\_\_\_\_。\underline{\hspace{2cm}}

  \centerquestionimage{0.35\textwidth}{chuyi_xia_2-p036-cleaned-9.jpg}

\item
  如图，有八个全等的三角形拼成一个大四边形 ABCD 和中间一个小四边形 MNPQ，连接 $EF$、GH 得到四边形 EFGH，设 $S_{四边形ABCD}=S_{1}$， $S_{四边形EFGH}=S_{2}$， $S_{四边形MNPQ}=S_{3}$，若 $S_{1}+S_{2}+S_{3}=10$，则 $S_{2}=$ \_\_\_\_ .\underline{\hspace{2cm}}

\item
  如图，在 10 个边长都为 1 的小正三角形的网格中，点 $P$ 是网格的一个顶点，以点 $P$ 为顶点作格点平行四边形（即顶点均在洛点上的四边形），请你写出所有可能的平行四边形的对角线的长 \_\_\_\_\underline{\hspace{2cm}}

  \begin{questionimagegroup}
  \questionimageitem{0.32\textwidth}{chuyi_xia_2-p036-cleaned-10.jpg}{}
  \questionimagegap
  \questionimageitem{0.32\textwidth}{chuyi_xia_2-p036-cleaned-11.jpg}{}
  \end{questionimagegroup}

\end{enumerate}

\section*{p037}
\subsection*{三、解答题（共 1 小题）}

\solutionlist[16]
\item
  如图， $\triangle ABC$ 是等边三角形，点 $D$ 是边 $BC$ 上的一点，以 $AD$ 为边作等边 $\triangle ADE$，过点 $C$ 作 $CF \parallel DE$ 交 $AB$ 于点 $F$ .
(1) 若点$D$是 $BC$边的中点（如图\circled{1}），求证：EF = CD；
(2) 在 (1) 的条件下直接写出 $\triangle AEF$ 和 $\triangle ABC$ 的面积比;
（3）若点 $D$ 是 $BC$ 边上的任意一点（除 $B$、 $C$ 外如图\circled{2}），那么（1）中的结论是否仍然成立？若成立，请给出证明；若不成立，请说明理由.

  \begin{questionimagegroup}
  \questionimageitem{0.32\textwidth}{chuyi_xia_2-p037-cleaned-1.jpg}{图\circled{1}}
  \questionimagegap
  \questionimageitem{0.32\textwidth}{chuyi_xia_2-p037-cleaned-2.jpg}{图\circled{2}}
  \end{questionimagegroup}

\end{enumerate}

\section*{p038}
\subsection*{四、未分类题目（共 2 小题）}

\solutionlist[17]
\item
  在 $\square ABCD$ 中， $\angle BAD$ 的平分线交直线 $BC$ 于点 $E$，交直线 $DC$ 于点 $F$ .
(1) 在图\circled{1} 中证明 CE = CF;
2\` 若 $\angle ABC=90^{\circ}$ , $G$ 是 $EF$ 的中点 (如图\circled{2}), 直接写出 $\angle BDG$ 的度数;
(3) 若 $\angle ABC=120^{\circ}$ , $FG \parallel CE$ , $FG=CE$ , 分别连接 $DB$、 $DG$ (如图\circled{3}), 求 $\angle BDG$ 的度数.

  \begin{questionimagegroup}
  \questionimageitem{0.28\textwidth}{chuyi_xia_2-p038-cleaned-1.jpg}{图\circled{1}}
  \questionimagegap
  \questionimageitem{0.28\textwidth}{chuyi_xia_2-p038-cleaned-2.jpg}{图\circled{2}}
  \questionimagegap
  \questionimageitem{0.28\textwidth}{chuyi_xia_2-p038-cleaned-3.jpg}{图\circled{3}}
  \end{questionimagegroup}

\item
  在 $\triangle ABC$ 中，点$P$为BC的中点.
（1）如图\circled{1}，求证： $AP < \frac{1}{2}$ （$AB+AC)$
(2) 延长 $AB$到 D，使得 BD = AC，延长 $AC$到 E，使得 CE = AB，连接 $DE$.
\circled{1}如图\circled{2}，连接 $BE$，若 $\angle BAC=60^{\circ}$，请你探究线段$BE$与线段$AP$之间的数量关系．写出你的结论，并加以证明；
\circled{2}请在图\circled{3}中证明： $BC \geqslant \frac{1}{2} DE.$

  \begin{questionimagegroup}
  \questionimageitem{0.28\textwidth}{chuyi_xia_2-p038-cleaned-4.jpg}{}
  \questionimagegap
  \questionimageitem{0.28\textwidth}{chuyi_xia_2-p038-cleaned-5.jpg}{图\circled{2}}
  \questionimagegap
  \questionimageitem{0.28\textwidth}{chuyi_xia_2-p038-cleaned-6.jpg}{图\circled{3}}
  \end{questionimagegroup}

\end{enumerate}

\section*{p039}
\subsection*{一、选择题（共 7 小题）}

\questionlist
\item
  \begin{questionwithimage}{chuyi_xia_2-p039-cleaned-1.jpg}
  如图， $\triangle ABC$ 的周长为 26，点 $D, E$ 都在边 $BC$ 上， $\angle ABC$ 的平分线垂直于 $AE$，垂足为 $Q$， $\angle ACB$ 的平分线垂直于 $AD$，垂足为 $P$，若 $BC=10$，则 $PQ$ 的长为（\quad）
  \optionlist
  \item $\frac{3}{2}$
  \item $\frac{5}{2}$
  \item 3
  \item 4
  \end{enumerate}
  \end{questionwithimage}

\item
  \begin{questionwithimage}{chuyi_xia_2-p039-cleaned-2.jpg}
  如图，在 $\triangle ABC$ 中， $M$ 为 $BC$ 中点， $AN$ 平分 $\angle BAC$， $AN \perp BN$ 于 $N$，且 $AB=10$， $AC=16$，则 $MN$ 等于（\quad）
  \optionlist
  \item 2
  \item 2.5
  \item 3
  \item 3.5
  \end{enumerate}
  \end{questionwithimage}

\item
  $\triangle ABC$ 的三边长分别为 $a 、 b 、 c$，三条中位线组成第一个中点三角形，第一个中点三角形的三条中位线又组成第二个中点三角形，以此类推，求第2019个中点三角形的周长为（\quad）
  \optionlist
  \item $\frac{a+b+c}{2^{2018}}$
  \item $\frac{a+b+c}{2^{2019}}$
  \item $\frac{a+b+c}{2^{2020}}$
  \item $\frac{a+b+c}{2^{2021}}$
  \end{enumerate}

\item
  \begin{questionwithimage}{chuyi_xia_2-p039-cleaned-3.jpg}
  如图，在 $\triangle ABC$ 中， $CE$、 $CF$ 分别平分 $\angle ACB$ 和 $\angle ACD$， $AE \parallel CF$ $AF \parallel CE$，直线 $EF$ 分别交 $AB$、 $AC$ 于点 $M$、 $N$ .若 $BC=a$， $AC=b$， $AB=c$，且 $c > a > b$，则 $ME$ 的长为（\quad）
  \optionlist
  \item $\frac{c-a}{2}$
  \item $\frac{a-b}{2}$
  \item $\frac{c-b}{2}$
  \item $\frac{a+b-c}{2}$
  \end{enumerate}
  \end{questionwithimage}

\item
  \begin{questionwithimage}{chuyi_xia_2-p039-cleaned-4.jpg}
  如图， $\triangle ABC$ 的周长为 16， $G$、 $H$ 分别为 $AB$、 $AC$ 的中点，分别以 $AB$、 $AC$ 为斜边向外作 Rt $\triangle ADB$ 和 Rt $\triangle AEC$，连接 $DG$、 $GH$、 $EH$，则 $DG+GH+EH$ 的值为（\quad）
  \optionlist
  \item 6
  \item 7
  \item 8
  \item 9
  \end{enumerate}
  \end{questionwithimage}

\item
  \begin{questionwithimage}{chuyi_xia_2-p039-cleaned-5.jpg}
  已知：四边形 $ABCD$ 中， $AB=2$， $CD=3$， $M$、 $N$ 分别是 $AD$， $BC$ 的中点，则线段 $MN$ 的取值范围是（\quad）
  \optionlist
  \item $1 < MN < 5$
  \item $1 < MN \leqslant 5$
  \item $\frac{1}{2} < MN < \frac{5}{2}$
  \item $\frac{1}{2} < MN \leqslant \frac{5}{2}$
  \end{enumerate}
  \end{questionwithimage}

\item
  \begin{questionwithimage}{chuyi_xia_2-p039-cleaned-6.jpg}
  如图，在 $\triangle ABC$ 中， $AB=AC=5$， $BC=6$， $D$、 $E$ 分别是 $AB$、 $AC$ 的中点， $F$、 $G$ 为 $BC$ 上的两点， $FG=3$，线段 $DG$， $EF$ 的交点为 $O$，当线段 $FG$ 在线段 $BC$ 上移动时，三角形 $FGO$ 的面积与四边 $ADOE$ 的面积之和恒为定值，则这个定值是（\quad）
  \optionlist
  \item 15
  \item 12
  \item 9
  \item 6
  \end{enumerate}
  \end{questionwithimage}

\end{enumerate}

\section*{p040}
\subsection*{四、未分类题目（共 1 小题）}

\solutionlist[8]
\item
  \begin{questionwithimage}{chuyi_xia_2-p040-cleaned-1.jpg}
  如图所示， $\triangle A_{0}BC$ 为等腰三角形，在线段$BC$上， $A_{1}$ 为 BC 中点， $A_{3}$ 为 $A_{1}C$ 中点， $A_{5}$ 为 $A_{3}C$ 中点， $A_{7}$ 为 $A_{5}C$ 中点，…，在线段 $A_{0}C$ 上， $A_{2}$ 为 $A_{0}C$ 中点， $A_{4}$ 为 $A_{2}C$ 中点， $A_{6}$ 为 $A_{4}C$ 中点…，依此下去，若 BC=6，A0B=5，则线段 $A_{n}A_{n+1}$ （n>0）的长度不可能为（\quad）
  \optionlist
  \item $\frac{5}{8}$
  \item $\frac{1}{8}$
  \item $\frac{1}{96}$
  \item $\frac{5}{128}$
  \end{enumerate}
  \end{questionwithimage}

\end{enumerate}

\subsection*{二、填空题（共 6 小题）}

\blanklist[9]
\item
  在四边形 $ABCD$ 中，对角线 $AC \perp BD$ 且 $AC=6$、 $BD=8$， $E$、 $F$ 分别是边 $AB$、 $CD$ 的中点，则 $EF=$ \_\_\_\_.\underline{\hspace{2cm}}

  \centerquestionimage{0.35\textwidth}{chuyi_xia_2-p040-cleaned-2.jpg}

\item
  如图，在 $\triangle ABC$ 中， $CD$ 是高， $CE$ 是中线， $CE=CB$，点 $A$、 $D$ 关于点 $F$ 对称，过点 $F$ 作 $FG \parallel CD$，交 $AC$ 边于点 $G$，连接 $GE$。若 $AC=18$， $BC=12$，则 $\triangle CEG$ 的周长为\_\_\_\_\_\_。\underline{\hspace{2cm}}

  \centerquestionimage{0.35\textwidth}{chuyi_xia_2-p040-cleaned-3.jpg}

\item
  如图，有一块直角三角形的木板 $AOB$， $\angle O=90^\circ$， $OA=3$， $OB=4$，一只小蚂蚁在 $OA$ 边上爬行（可以与 $O$、 $A$ 重合），设其所处的位置 $C$ 到 $AB$ 的中点 $D$ 的距离为 $x$，则 $x$ 的取值范围是 \_\_\_\_。\underline{\hspace{2cm}}

  \centerquestionimage{0.35\textwidth}{chuyi_xia_2-p040-cleaned-4.jpg}

\item
  如图：已知 $AB=10$，点 $C$、 $D$ 在线段 $AB$ 上且 $AC=DB=1$； $P$ 是线段 $CD$ 上的动点，分别以 $AP$、 $PB$ 为边在线段 $AB$ 的同侧作等边 $\triangle AEP$ 和等边 $\triangle PFB$，连接 $EF$，设 $EF$ 的中点为 $G$；当点 $P$ 从点 $C$ 运动到点 $D$ 时，则点 $G$ 移动路径的长是 \_\_\_\_。\underline{\hspace{2cm}}

  \centerquestionimage{0.35\textwidth}{chuyi_xia_2-p040-cleaned-5.jpg}

\item
  四边形 $ABCD$ 中， $AD > BC$， $E$、 $F$ 分别是 $AB$、 $CD$ 的中点， $AD$、 $BC$ 的延长线分别与 $EF$ 的延长线交于 $H$、 $G$，则 $\angle AHE$ \_\_\_\_ $\angle BGE$ （填“＞”或“＝”或“＜”号）\underline{\hspace{2cm}}

  \centerquestionimage{0.35\textwidth}{chuyi_xia_2-p040-cleaned-6.jpg}

\item
  如图，在四边形 $ABCD$ 中，对角线 $AC$、 $BD$ 相交于点 $O$，且 $AC=BD$， $E$、 $F$ 分别是 $AB$、 $CD$ 的中点， $EF$ 分别交 $BD$、 $AC$ 于点 $G$、 $H$，若 $\angle OBC=55^\circ$， $\angle OCB=45^\circ$，则 $\angle OGH=\_$ \_\_\_\_.\underline{\hspace{2cm}}

  \centerquestionimage{0.35\textwidth}{chuyi_xia_2-p040-cleaned-7.jpg}

\end{enumerate}

\section*{p041}
\subsection*{三、解答题（共 3 小题）}

\solutionlist[15]
\item
  如图所示，在梯形 $ABCD$ 中， $AD \parallel BC$， $AD < BC$， $F$， $E$ 分别是对角线 $AC$， $BD$ 的中点。求证： $EF=\frac{1}{2}(BC-AD)$ .

  \centerquestionimage{0.35\textwidth}{chuyi_xia_2-p041-cleaned-1.jpg}

\item
  如图，已知 $AE$、 $BD$ 相交于点 $C$， $AC=AD$， $BC=BE$， $F$、 $G$、 $H$ 分别是 $DC$、 $CE$、 $AB$ 的中点。求证： $^{4}$ （1） $HF=HG$；（2） $\angle FHG=\angle DAC$。

  \centerquestionimage{0.35\textwidth}{chuyi_xia_2-p041-cleaned-2.jpg}

\item
  如图\circled{1}，在四边形ABCD中， $AB=CD$， $E$， $F$ 分别是BC，AD的中点，连接 $EF$ 并延长，分别与BA，CD的延长线交于点 $M$， $N$，则 $\angle BME=\angle CNE$ （不需证明）.小明的思路是：在图\circled{1}中，连接 $BD$，取 $BD$的中点 $H$，连接 $HE$，HF，根据三角形中位线定理和平行线性质，可证得 $\angle BME=\angle CNE$
问题：如图\circled{2}，在 $\triangle ABC$ 中， $AC > AB$， $D$ 点在 $AC$ 上， $AB=CD$， $E$， $F$ 分别是 $BC$， $AD$ 的中点，连接 $EF$ 并延长，与 $BA$ 的延长线交于点 $G$，若 $\angle EFC=60^{\circ}$，连接 $GD$，判断 $\triangle AGD$ 的形状并证明.

  \begin{questionimagegroup}
  \questionimageitem{0.28\textwidth}{chuyi_xia_2-p041-cleaned-3.jpg}{图\circled{1}}
  \questionimagegap
  \questionimageitem{0.28\textwidth}{chuyi_xia_2-p041-cleaned-4.jpg}{图\circled{2}}
  \questionimagegap
  \questionimageitem{0.28\textwidth}{chuyi_xia_2-p041-cleaned-5.jpg}{}
  \end{questionimagegroup}

\end{enumerate}

\section*{p042}
\subsection*{四、未分类题目（共 1 小题）}

\solutionlist[18]
\item
  （1）如图所示，BD，CE分别是$\triangle ABC$的外角平分线，过点$A$作 $AF \perp BD$， $AG \perp CE$，垂足分别为F，G，连接 $FG$，延长 $AF$，AG，与直线$BC$分别交于点$M$、N，那么线段$FG$与$\triangle ABC$的周长之间存在的数量关系是什么？
即： $FG=$ \_\_\_\_ (AB+BC+AC) (直接写出结果即可)
（2）如图，若 $BD$， $CE$ 分别是 $\triangle ABC$ 的内角平分线；其他条件不变，线段 $FG$ 与 $\triangle ABC$ 三边之间又有怎样的数量关系？请写出你的猜想，并给予证明.
（3）如图，若 $BD$ 为 $\triangle ABC$ 的内角平分线， $CE$ 为 $\triangle ABC$ 的外角平分线，其他条件不变，线段 $FG$ 与 $\triangle ABC$ 三边又有怎样的数量关系？直接写出你的猜想即可．不需要证明．答：线段 $FG$ 与 $\triangle ABC$ 三边之间数量关系是\_\_\_\_.

  \begin{questionimagegroup}
  \questionimageitem{0.28\textwidth}{chuyi_xia_2-p042-cleaned-1.jpg}{}
  \questionimagegap
  \questionimageitem{0.28\textwidth}{chuyi_xia_2-p042-cleaned-2.jpg}{}
  \questionimagegap
  \questionimageitem{0.28\textwidth}{chuyi_xia_2-p042-cleaned-3.jpg}{}
  \end{questionimagegroup}

\end{enumerate}

\section*{p043}
\subsection*{一、选择题（共 8 小题）}

\questionlist
\item
  \begin{questionwithimage}{chuyi_xia_2-p043-cleaned-1.jpg}
  如图，E，F分别是矩形A$BCD$边$AD$、BC上的点，且 $\triangle ABG$， $\triangle DCH$ 的面积分别为15和20，则图中阴影部分的面积为（\quad）
  \optionlist
  \item 15
  \item 20
  \item 35
  \item 40
  \end{enumerate}
  \end{questionwithimage}

\item
  矩形 $ABCD$ 中, 对角线 $AC$、 $BD$ 交于点 $O$ , $\angle BOC=2\angle DOC$ , 若 $AC=8cm$ , 则 $AD$ 长为（\quad）
  \optionlist
  \item $4cm$
  \item $8cm$
  \item $\sqrt{48}cm$
  \item $\sqrt{80}cm$
  \end{enumerate}

\item
  \begin{questionwithimage}{chuyi_xia_2-p043-cleaned-2.jpg}
  如图，矩形纸片 ABCD 的边$AD$=9，AB=3，将其折叠，使点$D$与点$B$重合，则折叠后 DE 的长与折痕 EF 的长分别为（\quad）
  \optionlist
  \item $4, \sqrt{10}$
  \item $5, \sqrt{10}$
  \item $4, 2\sqrt{3}$
  \item $5, 2\sqrt{3}$
  \end{enumerate}
  \end{questionwithimage}

\item
  \begin{questionwithimage}{chuyi_xia_2-p043-cleaned-3.jpg}
  已知如图，矩形 ABCD 中 AB=4cm，BC=3cm，点$P$是 AB 上除 A，B 外任一点，对角线 AC，BD 相交于点$O$，DP，CP 分别交 AC；BD 于点$E$，F 且 $\triangle ADE$ 和 BCF 的面积之和 $4cm^{2}$，则四边形 PEOF 的面积为（\quad）
  \optionlist
  \item $1cm^{2}$
  \item $1.5cm^{2}$
  \item $2cm^{2}$
  \item $2.5cm^{2}$
  \end{enumerate}
  \end{questionwithimage}

\item
  \begin{questionwithimage}{chuyi_xia_2-p043-cleaned-4.jpg}
  如图，在矩形 ABCD 中， $AD=\sqrt{2}AB$，$\angle$BAD 的平分线交 BC 于点$E$， $DH \perp AE$ 于点$H$，连接 $BH$并延长交 CD 于点$F$，连接 $DE$交 BF 于点$O$，下列结论：
\circled{1} $\angle AED=\angle CED$ ; \circled{2} OE = OD; \circled{3} BH = HF; \circled{4} BC - CF = 2HE; \circled{5} AB = HF,
其中正确的有（\quad）
  \optionlist
  \item 2个
  \item 3个
  \item 4个
  \item 5个
  \end{enumerate}
  \end{questionwithimage}

\item
  \begin{questionwithimage}{chuyi_xia_2-p043-cleaned-5.jpg}
  已知，如图，在矩形 ABCD 中，P 是边$AD$上的动点，PE 垂直 AC 于 E， $PF \perp BD$ 于 F，如果 AB = 3，AD = 4，那么（\quad）
  \optionlist
  \item $PE+PF=\frac{12}{5}$
  \item $\frac{12}{5} < PE+PF < \frac{13}{5}$
  \item $PE+PF=5$
  \item 3 < PE + PF < 4
  \end{enumerate}
  \end{questionwithimage}

\item
  \begin{questionwithimage}{chuyi_xia_2-p043-cleaned-6.jpg}
  如图所示的矩形是由六个正方形组成，其中最小的正方形的面积为 1，则此矩形的面积为（\quad）
  \optionlist
  \item 99
  \item 120
  \item 143
  \item 168
  \end{enumerate}
  \end{questionwithimage}

\item
  矩形的一个内角的平分线把长边分成 $4\mathrm{~cm}$ 和 $6\mathrm{~cm}$ 两部分，则其面积为（\quad）

  \optionlist
  \item $24\mathrm{~cm}^{2}$
  \item $40\mathrm{~cm}^{2}$
  \item $60\mathrm{~cm}^{2}$
  \item $40\mathrm{~cm}^{2}$ 或 $60\mathrm{~cm}^{2}$
  \end{enumerate}

\end{enumerate}

\section*{p044}
\subsection*{二、填空题（共 6 小题）}

\blanklist[9]
\item
  如图，四边形 ABCD 是矩形，点$E$在线段$CB$的延长线上，连接 $DE$交 AB 于点$F$， $\angle AED=2\angle CED$，点$G$是 DF 的中点，若 BE = 1，AG = 4，则 AB 的长为 \_\_\_\_.\underline{\hspace{2cm}}

  \centerquestionimage{0.35\textwidth}{chuyi_xia_2-p044-cleaned-1.jpg}

\item
  Rt$\triangle ABC$中，$\angle$BAC=90$^\circ$，AB=3，AC=4，P 为边$BC$上一动点，PE$\perp$AB 于 E，PF$\perp$AC 于 F，M 为 EF 中点，则 AM 的最小值为 \_\_\_\_。\underline{\hspace{2cm}}

  \centerquestionimage{0.35\textwidth}{chuyi_xia_2-p044-cleaned-2.jpg}

\item
  矩形 $ABCD$ 中， $1B=4$， $BC=4\sqrt{3}$，点 $E$ 是折线段 $A-D-C$ 上的一个动点（点 $E$ 与点 $A$ 不重合），点 $P$ 是点 $A$ 关于 $BE$ 的对称点，在点 $E$ 运动的过程中，使 $\triangle PCB$ 为等腰三角形的点 $E$ 的位置共有 \_\_\_\_ 个.\underline{\hspace{2cm}}

\item
  如图，将矩形 ABCD 在直线 l 上按顺时针方向无滑动翻滚，可依次得到矩形 $A_{1}B_{1}C_{1}D$，矩形 $A_{2}B_{2}C_{1}D_{1}$，矩形 $A_{3}B_{2}C_{2}D_{2}$，…，若 AB=1，BC=2，那么 $AA_{12}$ 的长为 \_\_\_\_.\underline{\hspace{2cm}}

  \centerquestionimage{0.35\textwidth}{chuyi_xia_2-p044-cleaned-3.jpg}

\item
  如图，在矩形 $ABCD$ 中， $AD=6$， $AB=4$，点 $E$、 $G$、 $H$、 $F$ 分别在 $AB$、 $BC$、 $CD$、 $AD$ 上，且 $AF=CG=2$， $BE=DH=1$，点 $P$ 是直线 $EF$、 $GH$ 之间任意一点，连接 $PE$、 $PF$、 $PG$、 $PH$，则 $\triangle PEF$ 和 $\triangle PGH$ 的面积和等于 \_\_\_\_.\underline{\hspace{2cm}}

  \centerquestionimage{0.35\textwidth}{chuyi_xia_2-p044-cleaned-4.jpg}

\item
  如图，点 $A, B, C$ 在同一直线上，且 $BC=2AB$，点 $D, E$ 分别是 $AB, BC$ 的中点，分别以 $AB, DE, BC$ 为边，在 $A, C$ 同侧作三个正方形，得到三个平行四边形（阴影部分）的面积分别记作 $S_1, S_2, S_3$，若 $S_1+S_3=5$，则 $S_2=$ \_\_\_\_.\underline{\hspace{2cm}}

  \centerquestionimage{0.35\textwidth}{chuyi_xia_2-p044-cleaned-5.jpg}

\end{enumerate}

\subsection*{三、解答题（共 1 小题）}

\solutionlist[15]
\item
  \begin{questionwithimage}{chuyi_xia_2-p044-cleaned-6.jpg}
  我们定义：
  如图\circled{1}，在矩形 $MNPQ$ 中，点 $K$、$O$、$G$、$H$ 分别在线段 $NP$、$PQ$、$QM$、$MN$ 上，若 $\angle 1=\angle 2=\angle 3=\angle 4$，则称四边形 $KOGH$ 为矩形 $MNPQ$ 的反射四边形。

  如图\circled{2}、图\circled{3}，四边形 $ABCD$ 与 $A'B'C'D'$ 均为矩形，它们都是由 32 个边长为 1 的正方形组成的图形。点 $E$、$F$ 分别在 $BC$、$CD$ 边上，点 $E'$、$F'$ 分别在 $B'C'$、$C'D'$ 边上。请利用正方形网格，在图\circled{2}、图\circled{3}中分别画出矩形 $ABCD$ 和矩形 $A'B'C'D'$ 的反射四边形 $EFGH$ 和 $E'F'G'H'$。

  % 图注：图\circled{1}
  \end{questionwithimage}

\end{enumerate}

\section*{p045}
\subsection*{四、未分类题目（共 2 小题）}

\solutionlist[16]
\item
  如图，已知在矩形 ABCD 中，AC 是对角线。点$P$为矩形外一点且满足 AP=PC， $AP \perp PC$。PC 交 AD 于点$N$，连接 $DP$，过点$P$作 $PM \perp PD$ 交 AD 于 M。
（1）若 $AP=\sqrt{5}$， $AB=\frac{1}{3} BC$，求矩形ABCD的面积；
(2) 若 CD = PM，求证： $AC=AP+PN$ .

  \centerquestionimage{0.35\textwidth}{chuyi_xia_2-p045-cleaned-1.jpg}

\item
  阅读以下短文，然后解决下列问题：
如果一个三角形和一个矩形满足条件：三角形的一边与矩形的一边重合，且三角形的这边所对的顶点在矩形这边的对边上，则称这样的矩形为三角形的“友好矩形”，如图\circled{1}所示，矩形ABEF即为$\triangle ABC$的“友好矩形”，显然，当$\triangle ABC$是钝角三角形时，其“友好矩形”只有一个.
(1) 仿照以上叙述, 说明什么是一个三角形的 “友好平行四边形”；
（2）如图\circled{2}，若 $\triangle ABC$ 为直角三角形，且 $\angle C=90^{\circ}$，在图\circled{2}中画出 $\triangle ABC$ 的所有“友好矩形”，并比较这些矩形面积的大小；
（3）若 $\triangle ABC$ 是锐角三角形，且 $BC > AC > AB$，在图\circled{3}中画出 $\triangle ABC$ 的所有“友好矩形”，指出其中周长最小的矩形并加以证明.

  \centerquestionimage{0.35\textwidth}{chuyi_xia_2-p045-cleaned-2.jpg}

\end{enumerate}

\section*{p046}
\subsection*{四、未分类题目（共 2 小题）}

\solutionlist[18]
\item
  如图\circled{1}, 矩形 ${ABCD}$ 中, ${BC}={2AB}$ , $M$ 为 ${AD}$ 的中点,连接 ${BM}$ .
（1）请你判断并写出 $\angle BMD$ 是 $\angle ABM$ 的几倍；
（2）如图\circled{2}，在$\square$ABCD中，BC=2AB，M为AD的中点，CE$\perp$AB，连接 $EM$、CM，请问：$\angle$AEM与$\angle$DME是否也具有（1）中的倍数关系？若有，请证明；若没有，请说明理由.

  \begin{questionimagegroup}
  \questionimageitem{0.32\textwidth}{chuyi_xia_2-p046-cleaned-1.jpg}{图\circled{1}}
  \questionimagegap
  \questionimageitem{0.32\textwidth}{chuyi_xia_2-p046-cleaned-2.jpg}{图\circled{2}}
  \end{questionimagegroup}

\item
  如图，在 $\square ABCD$ 中， $AE \perp BC$， $AF \perp CD$， $E$， $F$ 为垂足， $\triangle AEF$ 的高线 $AN$， $FM$ 相交于点 $H$，设 $EF=a$， $AH=b$ ($a > b$ )，求 $\square ABCD$ 的对角线 $AC$ 的长.

  \centerquestionimage{0.35\textwidth}{chuyi_xia_2-p046-cleaned-3.jpg}

\end{enumerate}

\section*{p047}
\subsection*{一、选择题（共 8 小题）}

\questionlist
\item
  如图，公路 $AC$， $BC$ 互相垂直，公路 $AB$ 的中点 $M$ 与点 $C$ 被湖隔开．若测得 $AM$ 的长为 $1.2km$，则 $M$， $C$ 两点间的距离为（\quad）
  \optionlist
  \item $0.5km$
  \item $0.6km$
  \item $0.9km$
  \item $1.2km$
  \end{enumerate}

\item
  \begin{questionwithimage}{chuyi_xia_2-p047-cleaned-1.jpg}
  如图，在 Rt$\triangle ABC$中， $\angle ACB=90^{\circ}$，CD 为 $AB$边上的高，若点$A$关于 CD 所在直线的对称点$E$恰好为 AB 的中点，则 $\angle B$ 的度数是（\quad）
  \optionlist
  \item $60^{\circ}$
  \item $45^{\circ}$
  \item $30^{\circ}$
  \item $75^{\circ}$
  \end{enumerate}
  \end{questionwithimage}

\item
  \begin{questionwithimage}{chuyi_xia_2-p047-cleaned-3.jpg}
  如图，$\square$ABCD 中，BC: AB=1:2，M 为 AB 的中点，连接 $MD$、MC，则 $\angle DMC$ 等于（\quad）
  \optionlist
  \item $30^{\circ}$
  \item $60^{\circ}$
  \item $90^{\circ}$
  \item $45^{\circ}$
  \end{enumerate}
  \end{questionwithimage}

\item
  直角三角形的周长为 $2+\sqrt{6}$ , 斜边上的中线长为 1, 则该三角形的面积等于 （\quad）
  \optionlist
  \item 1
  \item $\frac{1}{2}$
  \item $\frac{1}{4}$
  \item $\frac{3}{4}$
  \end{enumerate}

\item
  \begin{questionwithimage}{chuyi_xia_2-p047-cleaned-4.jpg}
  如图，在$\triangle ABC$中，$\angle$ACB=90$^\circ$，D在BC上，E是AB的中点，A（\quad）
  \optionlist
  \item CE相交于F，且AD=D
  \item 若$\angle$B=20$^\circ$，则$\angle$DFE等于（）
  \item $30^{\circ}$
  \item $40^{\circ}$
  \item $50^{\circ}$
  \item $60^{\circ}$
  \end{enumerate}
  \end{questionwithimage}

\item
  如图，在四边形 ABCD 中， $\angle ABC=\angle ADC=90^{\circ}$，点$E$、F 分别是对角线 A（\quad）
  \optionlist
  \item BD 的中点，则（）
  \item $EF \perp BD$
  \item $\angle AEF=\angle ABD$
  \item $EF=\frac{1}{2}(AB+CD)$
  \item $EF=\frac{1}{2}(CD-AB)$
  \end{enumerate}

\item
  如图，将 Rt$\triangle ABC$绕其顶点$C$顺时针旋转 $90^{\circ}$ 得到 $\triangle DEC$， $\angle ACB=90^{\circ}$，M、N 分别为 A（\quad）
  \optionlist
  \item DE 的中点，若 MN = 4，则 AB 的长为（）
  \item $4\sqrt{2}$
  \item 4
  \item $2\sqrt{2}$
  \item 8
  \end{enumerate}

\item
  \begin{questionwithimage}{chuyi_xia_2-p047-cleaned-5.jpg}
  如图， $\triangle ABC$ 中， $BC=18$，若 $BD \perp AC$ 于 $D$， $CE \perp AB$ 于 $E$， $F$、 $G$ 分别为 $BC$、 $DE$ 的中点，若 $ED=10$，则 $FG$ 的长为（\quad）
  \optionlist
  \item $2\sqrt{14}$
  \item 9
  \item 10
  \item 无法确定
  \end{enumerate}
  \end{questionwithimage}

\end{enumerate}

\section*{p048}
\subsection*{四、未分类题目（共 1 小题）}

\solutionlist[9]
\item
  \begin{questionwithimage}{chuyi_xia_2-p048-cleaned-1.jpg}
  如图，边长为 $a$ 的等边 $\triangle ABC$ 的顶点 $A, B$ 分别在 $x$ 轴正半轴和 $y$ 轴正半轴上运动，则动点 $C$ 到原点 $O$ 的距离的最大值是（\quad）
  \optionlist
  \item $\frac{\sqrt{3}-1}{2} a$
  \item $\frac{\sqrt{3}+1}{2} a$
  \item $\frac{\sqrt{6}-1}{2} a$
  \item $\frac{\sqrt{6}+1}{2} a$
  \end{enumerate}
  \end{questionwithimage}

\end{enumerate}

\subsection*{二、填空题（共 5 小题）}

\blanklist[10]
\item
  在 $\triangle ABC$ 中， $\angle ACB=90^{\circ}$， $M$ 是 $AB$ 的中点， $E$、 $F$ 分别是 $AC$、 $BC$ 延长线上的点，且 $CE=CF=\frac{1}{2}AB$，则 $\angle EMF$ 的度数为\_\_\_\_.\underline{\hspace{2cm}}

  \centerquestionimage{0.35\textwidth}{chuyi_xia_2-p048-cleaned-2.jpg}

\item
  把一副三角板如图放置，E 是 AB 的中点，连接 $CE$、DE、CD，F 是 CD 的中点，连接 $EF$。若 AB=4，则 $S_{\triangle CEF}=$ \_\_\_\_。\underline{\hspace{2cm}}

  \centerquestionimage{0.35\textwidth}{chuyi_xia_2-p048-cleaned-3.jpg}

\item
  如图，在 $\triangle ABC$ 中， $AB=AC=7$， $BC=6$， $AF \perp BC$ 于 $F$， $BE \perp AC$ 于 $E$， $D$ 是 $AB$ 的中点，则 $\triangle DEF$ 的周长是\_\_\_\_.\underline{\hspace{2cm}}

\item
  如图，在 $\triangle ABC$ 中， $\angle BAC=90^\circ$，点 $D$ 在 $BC$ 延长线上，且 $AD=\frac{1}{2} BC$，若 $\angle D=50^\circ$，则 $\angle B=$ \_\_\_\_.\underline{\hspace{2cm}}

  \begin{questionimagegroup}
  \questionimageitem{0.32\textwidth}{chuyi_xia_2-p048-cleaned-4.jpg}{}
  \questionimagegap
  \questionimageitem{0.32\textwidth}{chuyi_xia_2-p048-cleaned-5.jpg}{}
  \end{questionimagegroup}

\item
  如图，Rt$\triangle ABC$中， $\angle ACB=90^{\circ}$， $\angle A=30^{\circ}$，CD 是 $\angle ACB$ 的角平分线，MD 是 $AB$边的中垂线，则 $\angle CDM=$ \_\_\_\_ .\underline{\hspace{2cm}}

\end{enumerate}

\section*{p049}
\subsection*{四、未分类题目（共 3 小题）}

\solutionlist[15]
\item
  如图，在 Rt$\triangle ABC$中，$\angle$ACB=90$^\circ$，CD 平分$\angle$A $\sim$ CE 是 $AB$边上的中线，CF$\perp$AB。求证：CD 平分$\angle$ECF。

  \centerquestionimage{0.35\textwidth}{chuyi_xia_2-p049-cleaned-1.jpg}

\item
  如图， $\angle ABE=\angle BAD=90^\circ$， $\angle AED=2\angle BED$，点 $G$ 是 $DF$ 的中点，若 $BE=1, AG=4$ .
(1) 求证: AE = AG;
(2) 求 AB 的长度.

  \centerquestionimage{0.35\textwidth}{chuyi_xia_2-p049-cleaned-2.jpg}

\item
  如图，在图甲中，已知 $AB=BC$， $\angle ABC=90^\circ$，点 $D$ 在 $AC$ 上， $DE \perp AC$ 于 $D$，交 $AB$ 于点 $E$，点 $M$ 为 $EC$ 的中点.
(1) 猜想线段 $BM$ 与 $DM$ 之间有什么关系？写出你的猜想，并加以证明.
(2) 如图乙, 将 $\triangle ADE$ 绕点 $A$ 逆时针旋转 $180^{\circ}$ , 第 (1) 问中的结论是否仍然成立? 请说明理由.

  \centerquestionimage{0.35\textwidth}{chuyi_xia_2-p049-cleaned-3.jpg}

\end{enumerate}

\section*{p050}
\subsection*{四、未分类题目（共 2 小题）}

\solutionlist[18]
\item
  如图，四边形 $ABCD$ 是矩形，点 $E$ 在线段 $CB$ 的延长线上，连接 $DE$ 交 $AB$ 于点 $F$， $\angle AED=2\angle CED$，点 $G$ 是 $DF$ 的中点.
（1）若 $BE=2$， $AG=4$，求 $AB$ 的长；
(2) 若 BC=2AB，求 $\angle AED$ 的度数.

  \centerquestionimage{0.35\textwidth}{chuyi_xia_2-p050-cleaned-1.jpg}

\item
  有一个直角三角形纸片 $BCE$，设点 $A$ 是斜边 $BE$ 上的一点，连接 $AC$，现沿 $AC$ 将纸片剪开，并将纸片 $ADE$ 顺时针旋转摆放成图\circled{2}、图\circled{3}、或图\circled{4} 的样子.
(1) 如图\circled{2}, 当点 $A$ 是中点, 且 ${DE}//{BC}$ 时,求 $\angle {BAE}$ 的度数；
（2）如图\circled{3}，当点 $A$ 是中点，但 $DE$ 不平行于 $BC$ 时，设 $M$ 是 $DE$ 的中点，连接 $AM$ 交 $BC$ 于点 $N$，求证： $\angle ANB+\angle BAE=180^{\circ}$
（3）如图\circled{4}，当 $AB < AE$ 时，设 $M$ 是 $DE$ 上的一点，连接 $AM$ 交 $BC$ 于点 $N$，若 $\angle ANB+\angle BAE=180^{\circ}$，那么点 $M$ 在 $DE$ 上的位置满足什么条件？

  \begin{questionimagegroup}
  \questionimageitem{0.28\textwidth}{chuyi_xia_2-p050-cleaned-2.jpg}{}
  \questionimagegap
  \questionimageitem{0.28\textwidth}{chuyi_xia_2-p050-cleaned-3.jpg}{图\circled{2}}
  \questionimagegap
  \questionimageitem{0.28\textwidth}{chuyi_xia_2-p050-cleaned-4.jpg}{图\circled{3}}
  \questionimagegap
  \questionimageitem{0.28\textwidth}{chuyi_xia_2-p050-cleaned-5.jpg}{图\circled{4}}
  \end{questionimagegroup}

\end{enumerate}

\section*{p051}
\subsection*{一、选择题（共 8 小题）}

\questionlist
\item
  下列说法正确的是（\quad）
  \optionlist
  \item 两组对角分别相等的四边形是矩形
  \item 有两个角是直角的四边形是矩形
  \item 一个角是直角的平行四边形是矩形
  \item 有一个角是直角，有一组对边相等的四边形是矩形
  \end{enumerate}

\item
  已知 $\square ABCD$，对角线交于点 $O_3$ 下列条件中不一定能确定为矩形的是（\quad）
  \optionlist
  \item $\angle ABC=90^\circ$
  \item $OA=OB$
  \item $AB=BC$
  \item $AC=BD$
  \end{enumerate}

\item
  检查一个门框是矩形的方法是（\quad）
  \optionlist
  \item 测量两条对角线是否相等
  \item 测量有三个角是直角
  \item 测量两条对角线是否互相平分
  \item 测量两条对角线是否互相垂直
  \end{enumerate}

\item
  \begin{questionwithimage}{chuyi_xia_2-p051-cleaned-1.jpg}
  已知 $AC$ 为矩形 $ABCD$ 的对角线, 则图中 $\angle 1$ 与 $\angle 2$ 一定不相等的是 （\quad）
  \optionlist
  \item 
  \item 
  \item 
  \item 
  \end{enumerate}
  \end{questionwithimage}

\item
  \begin{questionwithimage}{chuyi_xia_2-p051-cleaned-5.jpg}
  如图, 有两张形状、大小完全相同的直角三角形纸片 (同一个直角三角形的两条直角边不相等), 把两个三角形相等的边靠在一起 (两张纸片不重叠), 可以拼出若干种图形, 其中, 形状不同的四边形有 （\quad）
  \optionlist
  \item 3 种
  \item 4 种
  \item 5 种
  \item 6 种
  \end{enumerate}
  \end{questionwithimage}

\item
  \begin{questionwithimage}{chuyi_xia_2-p051-cleaned-6.jpg}
  如图，四边形 ABCD 的对角线 A（\quad）
  \optionlist
  \item BD 相交于点$O$，已知下列 6 个条件：
\circled{1}AB // DC；\circled{2}AB = DC；\circled{3}AC = BD；\circled{4}$\angle$ABC = 90$^\circ$；\circled{5}OA = OC；\circled{6}OB = O
  \item 
  \end{enumerate}
  \end{questionwithimage}

\item
  \begin{questionwithimage}{chuyi_xia_2-p051-cleaned-7.jpg}
  如图，E，F，G，H分别是四边形ABCD四条边的中点，要使四边形EFGH为矩形，则四边形ABCD应具备的条件是（\quad）
  \optionlist
  \item 一组对边平行而另一组对边不平行
  \item 对角线相等
  \item 对角线互相垂直
  \item 对角线互相平分
  \end{enumerate}
  \end{questionwithimage}

\item
  \begin{questionwithimage}{chuyi_xia_2-p051-cleaned-8.jpg}
  如图，过四边形 ABCD 的四个顶点分别作对角线 A（\quad）
  \optionlist
  \item BD 的平行线，若所围成的四边形 EFGH 是矩形时，原四边形 ABCD 必须满足的条件是（）
  \item $AD \perp CD$
  \item AD = CD
  \item $AC \perp BD$
  \item AC = BD
  \end{enumerate}
  \end{questionwithimage}

\end{enumerate}

\section*{p052}
\subsection*{二、填空题（共 6 小题）}

\blanklist[9]
\item
  如图，矩形 ABCD 中，点$A$（-4，1）、B（0，1）、C（0，3），则点$A$到 x 轴的距离是 \_\_\_\_，点$A$关于 x 轴的对称点$A$' 坐标是 \_\_\_\_，点$D$坐标是 \_\_\_\_。\underline{\hspace{2cm}}

  \centerquestionimage{0.35\textwidth}{chuyi_xia_2-p052-cleaned-1.jpg}

\item
  在四边形 $ABCD$ 中， $AB=DC$， $AD=BC$，请再添加一个条件，使四边形 $ABCD$ 是矩形。你添加的条件是 \_\_\_\_。（写出一种即可）\underline{\hspace{2cm}}

\item
  如图，矩形 OABC 的顶点$A$、C 在坐标轴上，B（4，-2），经过定点$P$的直线 l 把矩形 OABC 分成面积相等的两部分，则 P 点坐标为 \_\_\_\_。\underline{\hspace{2cm}}

  \centerquestionimage{0.35\textwidth}{chuyi_xia_2-p052-cleaned-2.jpg}

\item
  用一刻度尺检验一个四边形是否为矩形, 以下方法可行的有\_\_\_\_\_. (只要填序号即可)
\circled{1}量出四边及两条对角线，比较对边是否相等，对角线是否相等.
\circled{2}量出对角线的交点到四个顶点的距离，看是否相等.
\circled{3}量出一组邻边的长 $a$、 $b$ 以及和这两边组成三角形的那条对角线的长 $c$，计算是否有 $a^2+b^2=c^2$
\circled{4}量出两条对角线长，看是否相等.\underline{\hspace{2cm}}

\item
  在平面直角坐标系上，有点 $A(-2,-2)$、$B(2, 2)$、$C(0, 4)$，当点 $D$ 的坐标为 \_\_\_\_ 时，四边形 $ABCD$ 是矩形。\underline{\hspace{2cm}}

\item
  如图，已知平行四边形 $ABCD$ 的对角线 $AC$、$BD$ 交于点 $O$，且 $DE \parallel AC$、$CE \parallel BD$。要使四边形 $OCED$ 是矩形，则平行四边形 $ABCD$ 还必须添加的条件是 \_\_\_\_（填一个即可）。\underline{\hspace{2cm}}

  \centerquestionimage{0.35\textwidth}{chuyi_xia_2-p052-cleaned-3.jpg}

\end{enumerate}

\subsection*{三、解答题（共 1 小题）}

\solutionlist[15]
\item
  在 $\square ABCD$ 中，对角线 $AC$、$BD$ 相交于点 $O$，直线 $EF$ 过点 $O$，且 $AF \perp BC$。求证：四边形 $AFCE$ 是矩形。

  \centerquestionimage{0.35\textwidth}{chuyi_xia_2-p052-cleaned-4.jpg}

\end{enumerate}

\section*{p053}
\subsection*{四、未分类题目（共 3 小题）}

\solutionlist[16]
\item
  \begin{questionwithimage}{chuyi_xia_2-p053-cleaned-1.jpg}
  如图，直角坐标系中的网格由单位正方形构成， $\triangle ABC$ 中， $A$ 点坐标为（2，3）
（1）若以（\quad）
  \optionlist
  \item 
  \item C 及点$D$的顶点的四边形为矩形，直接写出 D 点坐标 \_\_\_\_。
(2) 若以
  \item 
  \item C 及点$E$为顶点的四边形为平行四边形，试在图中画出所有 E 点的位置．并求出这些平行四边形对角线的最大值和最小值.
  \end{enumerate}
  \end{questionwithimage}

\item
  如图，已知 $\square ABCD$ 的两条对角线相交于点 $O$， $BE \perp AC$， $CF \perp BD$，垂足分别为 $E$， $F$。且 $BE=CF$。求证： $\square ABCD$ 是矩形。

  \centerquestionimage{0.35\textwidth}{chuyi_xia_2-p053-cleaned-2.jpg}

\item
  如图，在 $\triangle ABC$ 中，点 $O$ 是 $AC$ 边上（端点除外）的一个动点，过点 $O$ 作直线 $MN \parallel BC$。设 $MN$ 交 $\angle BCA$ 的平分线于点 $E$，交 $\angle BCA$ 的外角平分线于点 $F$，连接 $AE$、 $AF$。那么当点 $O$ 运动到何处时，四边形 $AECF$ 是矩形？并证明你的结论。

  \centerquestionimage{0.35\textwidth}{chuyi_xia_2-p053-cleaned-3.jpg}

\end{enumerate}

\section*{p054}
\subsection*{四、未分类题目（共 1 小题）}

\solutionlist[19]
\item
  如图，在 $\square ABCD$ 中，对角线 $BD=12cm$， $AC=16cm$， $AC$， $BD$ 相交于点 $O$，若 $E$， $F$ 是 $AC$ 上两动点，分别从 $A$， $C$ 两点以相同的速度向 $C$、 $A$ 运动，其速度为 $0.5\mathrm{cm/s}$ .
（1）当 $E$ 与 $F$ 不重合时，四边形DEBF是平行四边形吗？说明理由；
（2）点 $E, F$ 在 $AC$ 上运动过程中，以 $D, E, B, F$ 为顶点的四边形是否可能为矩形？如能，求出此时的运动时间 $t$ 的值；如不能，请说明理由.

  \centerquestionimage{0.35\textwidth}{chuyi_xia_2-p054-cleaned-1.jpg}

\end{enumerate}

\section*{p055}
\subsection*{一、选择题（共 8 小题）}

\questionlist
\item
  \begin{questionwithimage}{chuyi_xia_2-p055-cleaned-1.jpg}
  如图，在四边形 $ABCD$ 中， $AB \perp BC$， $AB \parallel DC$， $AB$， $BC$， $CD$ 分别为 2，2，
$2\sqrt{3}+2$，则 $\angle BAD$ 的度数等于（\quad）
  \optionlist
  \item $120^{\circ}$
  \item $135^{\circ}$
  \item $150^{\circ}$
  \item 以上都不对
  \end{enumerate}
  \end{questionwithimage}

\item
  在四边形 $ABCD$ 中, $\angle A=60^{\circ}$ , $AB \perp BC$ , $CD \perp AD$ , $AB=4cm$ , $CD=2cm$ , 求四边形 $ABCD$ 的周长 （\quad）
  \optionlist
  \item $10+2\sqrt{3}$
  \item $8+2\sqrt{5}$
  \item $8+3\sqrt{5}$
  \item $10+2\sqrt{5}$
  \end{enumerate}

\item
  \begin{questionwithimage}{chuyi_xia_2-p055-cleaned-2.jpg}
  如图，在矩形 $COED$ 中，点 $D$ 的坐标是（1，3），则 $CE$ 的长是（\quad）
  \optionlist
  \item 3
  \item $2\sqrt{2}$
  \item $\sqrt{10}$
  \item 4
  \end{enumerate}
  \end{questionwithimage}

\item
  在四边形 $ABCD$ 中， $\angle A=60^{\circ}$， $\angle ABC=\angle ADC=90^{\circ}$， $BC=2$， $CD=11$，自 $D$ 作 $DH \perp AB$ 于 $H$，则 $DH$ 的长是（\quad）
  \optionlist
  \item 7.5
  \item 7
  \item 6.5
  \item 5.5
  \end{enumerate}

\item
  下列关于矩形的说法，正确的是（\quad）
  \optionlist
  \item 对角线相等的四边形是矩形
  \item 对角线互相平分的四边形是矩形
  \item 矩形的对角线互相垂直且平分
  \item 矩形的对角线相等且互相平分
  \end{enumerate}

\item
  \begin{questionwithimage}{chuyi_xia_2-p055-cleaned-3.jpg}
  如图， $\triangle ABC$ 中， $AC$ 的垂直平分线分别交 $AC$、 $AB$ 于点 $D$、 $F$， $BE \perp DF$ 交 $DF$ 的延长线于点 $E$，已知 $\angle A=30^{\circ}$， $BC=2$， $AF=BF$，则四边形 $BCDE$ 的面积是（\quad）
  \optionlist
  \item $2\sqrt{3}$
  \item $3\sqrt{3}$
  \item 4
  \item $4\sqrt{3}$
  \end{enumerate}
  \end{questionwithimage}

\item
  已知下列命题中:
（1）矩形是轴对称图形，且有两条对称轴；（2）两条对角线相等的四边形是矩形；
（3）有两个角相等的平行四边形是矩形；（4）两条对角线相等且互相平分的四边形是矩形。其中正确的有（\quad）
  \optionlist
  \item 4个
  \item 3个
  \item 2个
  \item 1个
  \end{enumerate}

\item
  \begin{questionwithimage}{chuyi_xia_2-p055-cleaned-4.jpg}
  如图，在 $\triangle ABC$ 中， $AB=3$， $AC=4$， $BC=5$， $P$ 为边 $BC$ 上一动点， $PE \perp AB$ 于 $E$， $PF \perp AC$ 于 $F$，则 $EF$ 的最小值为（\quad）
  \optionlist
  \item 2
  \item 2.2
  \item 8.4
  \item 2.5
  \end{enumerate}
  \end{questionwithimage}

\end{enumerate}

\section*{p056}
\subsection*{四、未分类题目（共 1 小题）}

\solutionlist[9]
\item
  \begin{questionwithimage}{chuyi_xia_2-p056-cleaned-1.jpg}
  如图所示，在四边形 $ABCD$ 中， $AB \perp BC$， $AB \perp AD$， $BD=BC$， $\angle C=60^\circ$，如果 $\triangle DBC$ 的周长为 $m$，则 $AD$ 的长为（\quad）
  \optionlist
  \item $\frac{m}{3}$
  \item $\frac{m}{6}$
  \item $\frac{m}{8}$
  \item $\frac{m}{12}$
  \end{enumerate}
  \end{questionwithimage}

\end{enumerate}

\subsection*{二、填空题（共 5 小题）}

\blanklist[10]
\item
  过 Rt$\triangle ABC$ 的斜边 $AB$ 上一点 $D$，作 $DE \perp AC$ 于点 $E$，$DF \perp BC$ 于点 $F$，则 $\angle FDE=$ \_\_\_\_。\underline{\hspace{2cm}}

\item
  如图，在 $\triangle ABC$ 中，$BC=8$，$AC=6$，$AB=10$，它们的中点分别是点 $D$、$E$、$F$，则 $CF$ 的长为 \_\_\_\_。\underline{\hspace{2cm}}

  \centerquestionimage{0.35\textwidth}{chuyi_xia_2-p056-cleaned-2.jpg}

\item
  如果四边形 $ABCD$ 的对角线 $AC, BD$ 相等，且互相平分于点 $O$，则四边形 $ABCD$ 是\_\_\_\_\_\_形，如果 $\angle AOB=60^\circ$，则 $AB: AC=$ \_\_\_\_.\underline{\hspace{2cm}}

\item
  在矩形 ABCD 中，M 为 $AD$边的中点，P 为 BC 上一点，PE$\perp$MC，PF$\perp$MB，当 AB、BC 满足条件 \_\_\_\_ 时，四边形 PEMF 为矩形.\underline{\hspace{2cm}}

  \centerquestionimage{0.35\textwidth}{chuyi_xia_2-p056-cleaned-3.jpg}

\item
  如图，在矩形ABCD中，BC=20cm，点$P$和点$Q$分别从点$B$和点$D$出发，按逆时针方向沿矩形ABCD的边运动，点$P$和点$Q$的速度分别为3cm/s和2cm/s，则最快\_\_\_\_s后，四边形ABPQ成为矩形.\underline{\hspace{2cm}}

  \centerquestionimage{0.35\textwidth}{chuyi_xia_2-p056-cleaned-4.jpg}

\end{enumerate}

\subsection*{三、解答题（共 1 小题）}

\solutionlist[15]
\item
  如图，在直角坐标系中， $\square ABCD$ 的顶点 $A, B, D$ 的坐标是 $A(-1, 0), B(0,-\frac{1}{2}), D(0, 2)$ .
(1) 求点$C$ 的坐标;
(2) 判断 $\square$ABCD 是不是矩形，并给出证明.

  \centerquestionimage{0.35\textwidth}{chuyi_xia_2-p056-cleaned-5.jpg}

\end{enumerate}

\section*{p057}
\subsection*{四、未分类题目（共 3 小题）}

\solutionlist[16]
\item
  如图，在矩形 $ABCD$ 中，对角线 $AC, BD$ 相交于点 $O$，点 $E, F, G, H$ 分别是 $AD, OA, BC, OC$ 的中点.
(1) 求证: 四边形 EFGH 为平行四边形.
(2) 当 AB, BC 满足什么条件时，四边形 EFGH 为矩形？并给出证明.

  \centerquestionimage{0.35\textwidth}{chuyi_xia_2-p057-cleaned-1.jpg}

\item
  如图, 在四边形 $ABCD$ 中, $AB \parallel CD$ , $AB=CD$ , $\angle A=\angle ADC$ , $E$ , $F$ 分别为 $AD$ , $CD$ 的中点, 连接 $BE$ , $BF$ , 延长 $BE$ 交 $CD$ 的延长线于点 $M$ .
(1) 求证: 四边形 $ABCD$ 为矩形;
(2) 若 MD=6, BC=12, 求 BF 的长度.

  \centerquestionimage{0.35\textwidth}{chuyi_xia_2-p057-cleaned-2.jpg}

\item
  如图， $O$ 是矩形 $ABCD$ 的对角线的交点， $E$、 $F$、 $G$、 $H$ 分别是 $OA$、 $OB$、 $OC$、 $OD$ 上的点，且 $AE=BF=CG=DH$ .
(1) 求证: 四边形 EFGH 是矩形;
（2）若 $E$、 $F$、 $G$、 $H$ 分别是 $OA$、 $OB$、 $OC$、 $OD$ 的中点，且 $DG \perp AC$， $OF=2cm$，求矩形ABCD的面积.

  \centerquestionimage{0.35\textwidth}{chuyi_xia_2-p057-cleaned-3.jpg}

\end{enumerate}

\section*{p058}
\subsection*{四、未分类题目（共 1 小题）}

\solutionlist[19]
\item
  已知矩形 $ABCD$ 和点 $P$ , 当点 $P$ 在 $BC$ 上任一位置（如图\circled{1}所示）时, 易证得结论: $PA^2+PC^2=PB^2+PD^2$ , 请你探究: 当点 $P$ 分别在图\circled{2}、图\circled{3}中的位置时, $PA^2$、 $PB^2$、 $PC^2$ 和 $PD^2$ 又有怎样的数量关系请你写出对上述两种情况的探究结论, 并利用图\circled{2}证明你的结论.
答: 对图\circled{2} 的探究结论为\_\_\_\_；
对图\circled{3}的探究结论为\_\_\_\_。

  \begin{questionimagegroup}
  \questionimageitem{0.28\textwidth}{chuyi_xia_2-p058-cleaned-1.jpg}{图\circled{1}}
  \questionimagegap
  \questionimageitem{0.28\textwidth}{chuyi_xia_2-p058-cleaned-2.jpg}{图\circled{2}}
  \questionimagegap
  \questionimageitem{0.28\textwidth}{chuyi_xia_2-p058-cleaned-3.jpg}{图\circled{3}}
  \end{questionimagegroup}

\end{enumerate}

\section*{p059}
\subsection*{一、选择题（共 7 小题）}

\questionlist
\item
  \begin{questionwithimage}{chuyi_xia_2-p059-cleaned-1.jpg}
  如图，四边形 ABCD 是菱形，以对角线 AC 为边向上作等边 $\triangle ACE$。已知 $\angle DAB=70^{\circ}$，则 $\angle EAD$ 的大小为（\quad）
  \optionlist
  \item $25^{\circ}$
  \item $35^{\circ}$
  \item $45^{\circ}$
  \item $55^{\circ}$
  \end{enumerate}
  \end{questionwithimage}

\item
  在菱形 $ABCD$ 中， $E, F$ 分别是 $BC, CD$ 上的点，若 $\triangle AEF$ 是等边三角形，且 $EF=AB$，则 $\angle BAD$ 的度数是（\quad）
  \optionlist
  \item $100^{\circ}$
  \item $105^{\circ}$
  \item $110^{\circ}$
  \item $120^{\circ}$
  \end{enumerate}

\item
  已知菱形的周长为 $8\sqrt{5}$ , 面积为 16, 则这个菱形较短的对角线长为 （\quad）
  \optionlist
  \item 4
  \item 8
  \item $4\sqrt{5}$
  \item 10
  \end{enumerate}

\item
  若菱形的周长为8，高为1，则菱形两邻角的度数比为（\quad）
  \optionlist
  \item 3:1
  \item 4:1
  \item 5:1
  \item 6:1
  \end{enumerate}

\item
  \begin{questionwithimage}{chuyi_xia_2-p059-cleaned-2.jpg}
  菱形 OABC 在平面直角坐标系中的位置如图所示，若 OA=2， $\angle AOC=60^{\circ}$，则 B 点的坐标是（\quad）
  \optionlist
  \item $(-1, \sqrt{3})$
  \item $(-3, \sqrt{3})$
  \item $(-3, \sqrt{2})$
  \item $(3, \sqrt{3})$
  \end{enumerate}
  \end{questionwithimage}

\item
  \begin{questionwithimage}{chuyi_xia_2-p059-cleaned-3.jpg}
  如图，P 为菱形 ABCD 的对角线 AC 上一点，AB=2cm， $\angle B=120^{\circ}$， $PE \perp CD$ 于点$E$， $PF \perp AD$ 于点$F$，则 $PE+PF$ 的值为（\quad）
  \optionlist
  \item $\sqrt{2}$
  \item $\sqrt{3}$
  \item 2
  \item $\frac{3\sqrt{2}}{2}$
  \end{enumerate}
  \end{questionwithimage}

\item
  \begin{questionwithimage}{chuyi_xia_2-p059-cleaned-4.jpg}
  如图，在菱形 ABCD 中， $\angle A=100^{\circ}$，E，F 分别是边$AB$和 BC 的中点， $EP \perp CD$ 于点$P$，则 $\angle FPC=（\quad）$
  \optionlist
  \item $35^{\circ}$
  \item $45^{\circ}$
  \item $50^{\circ}$
  \item $55^{\circ}$
  \end{enumerate}
  \end{questionwithimage}

\end{enumerate}

\section*{p060}
\subsection*{四、未分类题目（共 1 小题）}

\solutionlist[8]
\item
  \begin{questionwithimage}{chuyi_xia_2-p060-cleaned-1.jpg}
  如图，菱形 ABCD 中，AB=3，DF=1， $\angle DAB=60^{\circ}$， $\angle EFG=15^{\circ}$， $FG \perp BC$，则 AE=（\quad）
  \optionlist
  \item $1+\sqrt{2}$
  \item $\sqrt{6}$
  \item $2\sqrt{3}-1$
  \item $1+\sqrt{3}$
  \end{enumerate}
  \end{questionwithimage}

\end{enumerate}

\subsection*{二、填空题（共 6 小题）}

\blanklist[9]
\item
  在菱形 $ABCD$ 中，对角线 $BD, AC$ 交于点 $O, AE \perp BC$，且 $AE=OB$，则 $\angle CAE=\_\_\_\_$ .\underline{\hspace{2cm}}

  \centerquestionimage{0.35\textwidth}{chuyi_xia_2-p060-cleaned-2.jpg}

\item
  菱形的边长是 $2\mathrm{~cm}$ , 一条对角线的长是 $2\sqrt{3} \mathrm{~cm}$ , 则另一条对角线的长是 \_\_\_\_ cm.\underline{\hspace{2cm}}

\item
  菱形 $ABCD$ 中， $\angle A=60^{\circ}$， $AB=9$，点 $P$ 是菱形 $ABCD$ 内一点， $PB=PD=3\sqrt{3}$，则 $AP$ 的长为 \_\_\_\_.\underline{\hspace{2cm}}

\item
  已知菱形的边长为 $m$ ,对角线之和为 ${2n}$ ,则它的面积为\_\_\_\_.\underline{\hspace{2cm}}

\item
  如图，四边形 ABCD 是菱形， $\angle ABC=\angle \cap^{\circ}$，M 是 BC 的中点，CM = 2。点$P$是 BD 上一动点，则 $PM+PC$ 的最小值 \_\_\_\_.\underline{\hspace{2cm}}

  \centerquestionimage{0.35\textwidth}{chuyi_xia_2-p060-cleaned-3.jpg}

\item
  菱形 $ABCD, \angle ABC=60^{\circ}, B\Omega=6\sqrt{3}$ , 点 $E$ 在 $AB$ 上, $CE=2\sqrt{7}$ , 将 $CE$ 绕点 $C$ 旋转 $60^{\circ}$ 得到线段 $CF$ 交 $BD$ 于 $F$ , 则 $DF$ 的长为 \_\_\_\_\underline{\hspace{2cm}}

\end{enumerate}

\subsection*{三、解答题（共 1 小题）}

\solutionlist[15]
\item
  如图，四边形 $ABCD$ 是菱形，过 $AB$ 的中点 $E$ 作 $AC$ 的垂线 $EF$ 交 $AD$ 于 $M$，交 $CD$ 的延长线于点 $F$ .
（1）求证： $EF$、 $AD$ 互相平分；
(2) 若 DF=2，求菱形 ABCD 的周长.

  \centerquestionimage{0.35\textwidth}{chuyi_xia_2-p060-cleaned-4.jpg}

\end{enumerate}
